\documentclass[11pt]{article}
\usepackage[top=1.5cm, bottom=2cm, left=1.5cm, right=1.5cm]{geometry}
\usepackage{color}
%\usepackage{amsmath}

%%%
%\textheight 240mm % A4 297mm
%\textwidth 160mm   % A4 210mm
%\columnsep 10pt 
%\columnseprule 0pt 
%
%\oddsidemargin 0pt 
%\evensidemargin 0pt 
%\marginparwidth 90pt 
%\marginparsep 11pt 
%
%\topmargin 0pt 
%\headheight 0pt 
%\headsep 0pt 
%\footskip 30pt 
%
%\parskip 0pt plus 1pt 
%\parindent 1.5em 
%\partopsep 3pt plus 2pt minus 2pt 
%%%

\makeatletter
\def\eqnarray{\stepcounter{equation}\let\@currentlabel=\theequation
\global\@eqnswtrue
\global\@eqcnt\z@\tabskip\@centering\let\\=\@eqncr
$$\halign to \displaywidth\bgroup\@eqnsel\hskip\@centering
%%  $\displaystyle\tabskip\z@{##}$&\global\@eqcnt\@ne 
    $\displaystyle\tabskip\z@{##}$&\global\@eqcnt\@ne 
%%  \hskip 2\arraycolsep \hfil${##}$\hfil
                         \hfil$\displaystyle{{}##{}}$\hfil
%%  &\global\@eqcnt\tw@ \hskip 2\arraycolsep $\displaystyle\tabskip\z@{##}$\hfil 
    &\global\@eqcnt\tw@                      $\displaystyle\tabskip\z@{##}$\hfil 
   \tabskip\@centering&\llap{##}\tabskip\z@\cr}
\makeatother

\renewcommand{\labelenumi}{(\theenumi)}
\renewcommand{\theenumi}{\alph{enumi}}
\renewcommand{\theenumii}{(\arabic{enumi})}
\renewcommand{\theequation}{\arabic{section}.\arabic{equation}}

\def\baselinestretch{1.20}
\def\bvec#1{\leavevmode\hbox{\boldmath$#1$}}
\let\vec=\bvec
\def\r#1{{\rm#1}}
\def\ave#1{\left\langle#1\right\rangle}
\def\aves#1{\langle#1\rangle}
\def\dd#1#2{\frac{\partial #1}{\partial #2}}
\def\tensor#1{\mathord{\displaystyle\mathop{#1}^\leftrightarrow}}
\def\para{\parallel}
\def\ddrho{\frac{\partial}{\partial\rho}}
\def\ddV{\frac{\partial}{\partial V}}
\def\ddt{\frac{\partial}{\partial t}}
\def\ddr{\frac{\partial}{\partial r}}
\def\dds{\frac{\partial}{\partial s}}
\def\gradrho2{\aves{|\nabla\rho|^2}}
\def\psid{\dot{\psi}}
\def\psit{\psi_t}
\def\psitd{\dot{\psit}}
\def\Er{E_r}

\def\me{m_\r{e}}
\def\mi{m_\r{i}}
\def\mz{m_\r{z}}
\def\mb{m_\r{b}}
\def\nee{n_\r{e}}
\def\ni{n_\r{i}}
\def\nz{n_\r{z}}
\def\nb{n_\r{b}}
\def\Te{T_\r{e}}
\def\Ti{T_\r{i}}
\def\Tz{T_\r{z}}
\def\Tb{T_\r{b}}
\def\Zi{Z_\r{i}}
\def\Zz{Z_\r{z}}
\def\Zb{Z_\r{b}}
\def\uzt#1{u_{\r{#1}\zeta}}
\def\qzt#1{\hat{q}_{\r{#1}\zeta}}
\def\uszt{u_{s\zeta}}
\def\upara#1{u_{\r{#1}\para}}
\def\uhatth#1{\hat{u}_{\r{#1}\theta}}
\def\qhatpara#1{\hat{q}_{\r{#1}\para}}
\def\urho#1{u_\r{#1}^\rho}
\def\ugrho{u_{g}^\rho}
\def\uV#1{u_\r{#1}^V}
\def\ugV{u_{g}^V}
\def\chis#1{\chi_\r{#1}}
\def\De{D_\r{e}}
\def\nun#1{\nu_\r{0#1}}
\def\nsup#1{n_\r{#1sup}}
\def\Tsup#1{T_\r{#1sup}}
\def\nuLT#1{\nu_{\r{L}T_\r{#1}}}
\def\PRF#1{P_\r{RF#1}}
\def\Palp#1{P_{\alpha\r{#1}}}
\def\fei{f_\r{ei}}

\def\nbrp{n_\r{b}^\r{rp}}
\def\ubrp{u_\r{b}^\r{rp}}
\def\Dbrp{D_\r{b}^\r{rp}}
\def\Db{D_\r{b}}

\def\bri{\aves{B^2}\aves{R^2} - I^2}
%\def\rib{\aves{R^2}-\frac{I^2}{\aves{B^2}}}
\def\uspara{u_{s\para}}
\def\Pbpara{P_{\r{b}\parallel}}
\def\Pbperp{P_{\r{b}\perp}}
\def\PBr{P_\r{Br}}
\def\Pb{P_\r{b}}
\def\Eb{E_\r{b}}

\def\EH{E_\r{H}}
\def\nna{n_{01}}
\def\nnb{n_{02}}
\def\nnc{n_{03}}
\def\nnz{\bar{n}_{z}}
\def\Tna{T_{01}}
\def\Tnb{T_{02}}
\def\Tnc{T_{03}}

\def\nueff{\nu_\r{eff}}
\def\fCX{f_\r{CX}}
\def\nuL{\nu_\r{L}}
\def\nuCX{\nu_\r{CX}}
\def\nuLC{\nu_\r{LC}}
\def\nuion{\nu_\r{ionize}}
\def\nub{\nu_\r{b}}
\def\nuTei{\nu_\r{Tei}}
\def\vb{v_\r{b}}
\def\nuOL{\nu_\r{OL}}
\def\nubL{\nu_\r{bL}}

\newcommand{\Frac}[2]{%
  {\displaystyle {\displaystyle #1\over \displaystyle #2}}%
}
\newcommand{\rdprd}[2]{\Frac{\partial #1}{\partial #2}}
\newcommand{\aveaveB}[1]{\left< \!\! \left< #1 \right> \!\! \right>}
%\newcommand{\hatomega}{\widehat{\omega}}

\def\red#1{\textcolor{red}{#1}}
\def\blue#1{\textcolor{blue}{#1}}

\begin{document}

\begin{center}
TASK/TX V5.62 [26/09/29]
\end{center}

\section{Equations in $\rho$ coordinates}
\vspace{-5mm}

%
\begin{eqnarray}
%=======%
% \phi  %
%=======%
  0 &=&
%
    \frac{1}{V'} \ddrho \left[ V' \gradrho2 \ddrho{\Phi} \right]
%
  + \frac{1}{\epsilon_0}
%
    \left(-e \nee + \Zi e \ni + \Zz e \nz + \Zb e \nb + \Zb e g \nbrp \right)
\\[5mm]
%=========%
% \psitd  %
%=========%
  \frac{1}{c^2} \ddt \psitd &=&
%
    V' \ave{\frac{|\nabla\rho|^2}{R^2}} \ddrho \left[ \frac{1}{V'\aves{R^{-2}}}
				         \ddrho \psit \right] \nonumber\\[3mm]&&{}
%
  + \frac{\mu_0}{\aves{B^\theta}}
%
    \left[-     e \nee \left( \aves{B\upara{e}} - I\ave{\frac{\uzt{e}}{R}} \right)
          + \Zi e \ni  \left( \aves{B\upara{i}} - I\ave{\frac{\uzt{i}}{R}} \right)
    \right.\nonumber\\[3mm]&&{}
    \left.+ \Zz e \nz  \left( \aves{B\upara{z}} - I\ave{\frac{\uzt{z}}{R}} \right)
          + \Zb e \nb  \left( \aves{B\upara{b}} - I\ave{\frac{\uzt{b}}{R}} \right) \right]
\\[5mm]
%========%
% \psid  %
%========%
  \frac{1}{c^2} \ddt \psid &=&
%
  \frac{1}{V'\aves{R^{-2}}} \ddrho \left[ V'
				    \ave{\frac{|\nabla\rho|^2}{R^2}}
				    \ddrho \psi \right] \nonumber\\[3mm]&&{}
%
  - \frac{\mu_0}{\aves{R^{-2}}} \left(-e \nee \ave{\frac{\uzt{e}}{R}}
			         + \Zi e \ni  \ave{\frac{\uzt{i}}{R}}
			         + \Zz e \nz  \ave{\frac{\uzt{z}}{R}}
			         + \Zb e \nb  \ave{\frac{\uzt{b}}{R}} \right)
%
%  + \mu_0 j_\r{virtual}
\\[5mm]
%======%
% \psi %
%======%
  \ddt \psi &=&
%
  \psid
\\[5mm]
%========%
% \psit  %
%========%
  \ddt \psit &=&
%
  \psitd
\end{eqnarray}
%

%%%

%
\begin{eqnarray}
%=======%
% \nee  %
%=======%
  \frac{1}{V'} \ddt ( V' \nee ) &=&
%
  - \frac{1}{V'} \ddrho [ V' \nee (\urho{e} - \ugrho) ] 
%
  + \nuion\, \frac{\nee}{\nna + \nnb + \nnc}\, (\nna + \nnb + \nnc)
\nonumber\\[3mm]&&{}
%
  + \Zz \nee \nnz f_5S_5 - \Zz \nee \nz A_6
%
  - \nuL (\nee - \nsup{e})
%
  + (1 - \fCX) \frac{\Pb}{\Eb}
\\[5mm]
%=========%
% radial  %
%=========%
  0 &=&
%
  -        ( \bri ) \ddrho \nee \Te
%
  + e \nee ( \bri ) \ddrho \Phi
\nonumber\\[3mm]&&{}
%
  - e \dd{\psi}{\rho} I          \nee \aves{B\upara{e}}
%
  + e \dd{\psi}{\rho} \aves{B^2} \nee \aves{R\uzt{e}}
\\[5mm]
%===============%
% <B \upara{e}> %
%===============%
  \me \nee \ddt \aves{B \upara{e}} &=&
%
  - \hat{\mu}_1^\r{e} \aves{B \upara{e}} 
  - \hat{\mu}_2^\r{e} \aves{B \qhatpara{e}} 
%
  + \hat{\mu}_1^\r{e} {B V_{1\r{e}}} 
  + \hat{\mu}_2^\r{e} {B V_{2\r{e}}}
\nonumber\\[3mm]&&{}
%
  + \ell_{11}^{\r{ee}} \aves{B \upara{e}}
  + \ell_{11}^{\r{ei}} \aves{B \upara{i}}
  + \ell_{11}^{\r{ez}} \aves{B \upara{z}}
  + \ell_{11}^{\r{eb}} \aves{B \upara{b}}
\nonumber\\[3mm]&&{}
%
  - \ell_{12}^{\r{ee}} \aves{B \qhatpara{e}} 
  - \ell_{12}^{\r{ei}} \aves{B \qhatpara{i}} 
  - \ell_{12}^{\r{ez}} \aves{B \qhatpara{z}} 
\nonumber\\[3mm]&&{}
%
  - e \nee \aves{B E_\para}
%
%  + \aves{B S_{\r{m}\para}^\r{e}}
%
  - \nuL \me \nee \aves{B \upara{e}}
%
  - \nun{e} \me \nee \aves{B \upara{e}}
\nonumber\\[3mm]&&{}
%
  - \frac{1}{V'} \ddrho V' \aves{|\nabla \rho|} v_{\r{e}\zeta} \me \nee \aves{B \upara{e}}
%
  + \aves{|\nabla \rho|} v_{\r{e}\zeta} \me \nee 
  \biggl(\ddrho \frac{\aves{B^2}}{I} \biggl) \frac{I}{\aves{B^2}} \aves{B \upara{e}}
\nonumber\\[3mm]&&{}
%
  + \frac{1}{V'} \ddrho V' \aves{|\nabla \rho|^2} \chis{e\zeta} \me \nee
  \frac{\aves{B^2}}{I} \ddrho \frac{I}{\aves{B^2}} \aves{B \upara{e}}
\nonumber\\[3mm]&&{}
%
  - \aves{|\nabla \rho|^2} \chis{e\zeta} \me \nee \biggl( \ddrho
  \frac{\aves{B^2}}{I} \biggr) \ddrho \frac{I}{\aves{B^2}} \aves{B \upara{e}}
%
  - \frac{\aves{B^2}}{I} \frac{1}{V'} \ddrho V' \aves{\Pi_\r{e}^\r{res}}
\nonumber\\[3mm]&&{}
%
  - \frac{1}{V'} \ddrho V' 
    \left[  \aves{|\nabla\rho|} v_{\r{e}\zeta} \me \nee
          \left(\aves{R^2} - \frac{I^2}{\aves{B^2}}\right) \frac{1}{I}BV_{1\r{e}} \right.
\nonumber\\[3mm]&&{}
    \left.- \aves{|\nabla\rho|^2} \chis{e\zeta} \me \nee
          \ddrho \left(\aves{R^2} - \frac{I^2}{\aves{B^2}}\right) \frac{1}{I}BV_{1\r{e}}\right]
\\[5mm]
%================%
% elec. toroidal %
%================%
 \frac{1}{V'} \ddt (V' \me \nee \aves{R \uzt{e}}) &=&
%
  - \frac{1}{V'} \ddrho V' \left[  \aves{|\nabla\rho|} v_{\r{e}\zeta} \me \nee \aves{R \uzt{e}}
			       +   (\urho{e} - \ugrho) \me \nee \aves{R \uzt{e}} \right.
\nonumber\\[3mm]&&{}
                           \left.- \aves{|\nabla\rho|^2} \chis{e\zeta} \me \nee \ddrho \aves{R \uzt{e}}
                               + \aves{\Pi_\r{e}^\r{res}} \right]
\nonumber\\[3mm]&&{}
%
  + \ell_{11}^\r{ee} \aves{R \uzt{e}}
  + \ell_{11}^\r{ei} \aves{R \uzt{i}}
  + \ell_{11}^\r{ez} \aves{R \uzt{z}}
  + \ell_{11}^\r{eb} \aves{R \uzt{b}}
\nonumber\\[2.5mm]&&{}
%
  - \ell_{12}^\r{ee} \aves{R \qzt{e}}
  - \ell_{12}^\r{ei} \aves{R \qzt{i}}
  - \ell_{12}^\r{ez} \aves{R \qzt{z}}
\nonumber\\[2.5mm]&&{}
%
  - e                 \nee \aves{R E_\zeta}
%
  - e \dd{\psi}{\rho} \nee \urho{e}
%
\nonumber\\[2.5mm]&&{}
%
%  turbulent particle flux model
%
%  - \frac{e^2 \aves{|\nabla\rho|^2} \De}{\Te} \dd{\psi}{\rho} \nee
%    \left[  \dd{\psi}{\rho} \frac{\aves{B^2}}{\bri} \aves{R \uzt{e}}
%          - \dd{\psi}{\rho} \frac{I}{\bri} \aves{B \upara{e}} \right.
%\nonumber\\[2.5mm]&&{}
%    \left.+ \frac{C_\r{T}-1}{e} \ddrho \Te
%	  + \ddrho \Phi
%	  + \frac{\Te}{e} \dd{\psi}{\rho} C_\r{p} \right]
%
  - \frac{e \aves{|\nabla\rho|^2} \De}{\Te}  \dd{\psi}{\rho}
  \left[  (1-C_\r{T})\Te \ddrho \nee
        + C_\r{T} \ddrho \nee \Te
	+ \nee \Te \dd{\psi}{\rho} C_\r{p}\right]
\nonumber\\[2.5mm]&&{}
%
  - \nuL  \me \nee \aves{R \uzt{e}}
%
  - \nun{e} \me \nee \aves{R \uzt{e}}
\\[5mm]
%===========%
% \nee \Te  %
%===========%
 \frac{3}{2}\frac{1}{V'^{5/3}} \ddt (V'^{5/3} \nee \Te ) &=& 
%
  - \frac{1}{V'} \ddrho V'
    \left[ - \chis{e}       \aves{|\nabla\rho|^2} \nee \ddrho \Te
           + V_\r{e}^\r{hp} \aves{|\nabla\rho|}   \nee \Te
           + \frac{5}{2} \nee \Te ( \urho{e} - \ugrho ) \right]
\nonumber\\[2.5mm]&&{}
%
  - \ugrho \ddrho \nee \Te
%
  - \urho{i} \ddrho \ni  \Ti
  - \urho{z} \ddrho \nz  \Tz
  - \uhatth{i} \aves{\vec{B}\cdot\nabla\cdot\tensor{\pi}_\r{i}}
  - \uhatth{z} \aves{\vec{B}\cdot\nabla\cdot\tensor{\pi}_\r{z}}
\nonumber\\[2.5mm]&&{}
%
  - \frac{3}{2} \nee \frac{ \Te - \Ti }{\tau_\r{ei}}
  - \frac{3}{2} \nee \frac{ \Te - \Tz }{\tau_\r{ez}}
\nonumber\\[2.5mm]&&{}
%
%  - \dd{\rho}{\psi} \ddrho ( \nee \Te + \ni \Ti ) \psit
%  - \frac{1}{\mu_0} \dd{\rho}{\psi} \dd{I}{\rho} I \aves{R^{-2}}
%    \left( - \frac{\psitd}{q} + \psid \right)
%\nonumber\\[2.5mm]&&{}
%
  - \dot{\psi}_t \aves{B^\theta} (-     e \nee \uhatth{e}
                                  + \Zi e \ni  \uhatth{i}
                                  + \Zz e \nz  \uhatth{z})
\nonumber\\[2.5mm]&&{}
%  - \frac{1}{\aves{B^2}}\dot{\psi}_t (-e \nee \aves{Bu_{\r{e}\para}} +
%  Z_\r{i} e \ni \aves{Bu_{\r{i}\para}})
%
%  +\frac{1}{\aves{B^2}}\dot{\psi}_t (-e \nee BV_{1\r{e}} + Z_\r{i} e \ni
%  BV_{1\r{i}})
%\nonumber\\[2.5mm]&&{}
%
  +\dot{\psi}\left(-     e \nee \ave{\frac{\uzt{e}}{R}}
	           + \Zi e \ni  \ave{\frac{\uzt{i}}{R}}
	           + \Zz e \nz  \ave{\frac{\uzt{z}}{R}}\right)
\nonumber\\[2.5mm]&&{}
%
  - \nuL \Te (\nee - \nsup{e})
%
  - \frac{3}{2} \nuLT{e} \nee (\Te - \Tsup{e})
%
\nonumber\\[2.5mm]&&{}
  - \EH \nuion\, \frac{\nee}{\nna + \nnb + \nnc}\, (\nna + \nnb + \nnc)
\nonumber\\[2.5mm]&&{}
%
  + \fei \Pb
%
%  + \mb \frac{\Pb}{\Eb} \aves{B \upara{b}} \frac{\aves{B \upara{e}}}{\aves{B^2}}
%
  + \Palp{e}
%
  + \PRF{e}
%
  - \PBr
\\[5mm]
%=====================%
% 2<B \qepara>/(5 pe) %
%=====================%
 \frac{5}{2} \me \nee \ddt \aves{B \qhatpara{e}} &=&
%
  - \hat{\mu}_2^\r{e} \aves{B \upara{e}} 
  - \hat{\mu}_3^\r{e} \aves{B \qhatpara{e}}
%
  + \hat{\mu}_2^\r{e} {B V_{1\r{e}}} 
  + \hat{\mu}_3^\r{e} {B V_{2\r{e}}}
\nonumber\\[3mm]&&{}
%
  - \ell_{21}^{\r{ee}} \aves{B \upara{e}}
  - \ell_{21}^{\r{ei}} \aves{B \upara{i}}
  - \ell_{21}^{\r{ez}} \aves{B \upara{z}}
  - \ell_{21}^{\r{eb}} \aves{B \upara{b}}
\nonumber\\[3mm]&&{}
%
  + \ell_{22}^{\r{ee}} \aves{B \qhatpara{e}}
  + \ell_{22}^{\r{ei}} \aves{B \qhatpara{i}}
  + \ell_{22}^{\r{ez}} \aves{B \qhatpara{z}}
\nonumber\\[3mm]&&{}
%
  - \nuL \me \nee \aves{B \qhatpara{e}}
%
  - \nun{e} \me \nee \aves{B \qhatpara{e}}
\end{eqnarray}
%

%
\begin{eqnarray}
%======%
% \ni  %
%======%
  \frac{1}{V'} \ddt ( V' \ni ) &=&
%
  - \frac{1}{V'} \ddrho [ V' \ni (\urho{i} - \ugrho) ]
%
  + \nuion\, \frac{\nee}{\nna + \nnb + \nnc}\, (\nna + \nnb
  + \nnc)
\nonumber\\[2.5mm]&&{}
%
  - \nuL (\ni - \nsup{i}) 
%
  + \nub \nb
%
  + \nub g \nbrp
%
  - \fCX \frac{\Pb}{\Eb}
%
  + \nubL \nb
%
  + S_\r{LC}
%
  - \nuOL \ni
%
\\[5mm]
%=========%
% radial  %
%=========%
  0 &=&
%
  -           ( \bri ) \ddrho \ni \Ti
%
  - \Zi e \ni ( \bri ) \ddrho \Phi
\nonumber\\[3mm]&&{}
%
  + \Zi e \ni \dd{\psi}{\rho} I          \aves{B\upara{i}}
%
  - \Zi e \ni \dd{\psi}{\rho} \aves{B^2} \aves{R\uzt{i}}
\\[5mm]
%===============%
% <B \upara{i}> %
%===============%
  \mi \ni \ddt \aves{B \upara{i}} &=&
%
  - \hat{\mu}_1^\r{i} \aves{B \upara{i}} 
  - \hat{\mu}_2^\r{i} \aves{B \qhatpara{i}}
%
  + \hat{\mu}_1^\r{i} {B V_{1\r{i}}} 
  + \hat{\mu}_2^\r{i} {B V_{2\r{i}}}
\nonumber\\[3mm]&&{}
%
  + \ell_{11}^{\r{ie}} \aves{B \upara{e}}
  + \ell_{11}^{\r{ii}} \aves{B \upara{i}}
  + \ell_{11}^{\r{iz}} \aves{B \upara{z}}
  + \ell_{11}^{\r{ib}} \aves{B \upara{b}}
\nonumber\\[3mm]&&{}
%
  - \ell_{12}^{\r{ie}} \aves{B \qhatpara{e}}
  - \ell_{12}^{\r{ii}} \aves{B \qhatpara{i}}
  - \ell_{12}^{\r{iz}} \aves{B \qhatpara{z}}
\nonumber\\[3mm]&&{}
%
  + \Zi e \ni \aves{B E_\para}
%
%  + \aves{B S_{\r{m}\para}^\r{i}}
%
%  + \aves{B S_\r{LC\para}}
%
  - \nuL  \mi \ni \aves{B \upara{i}}
%
  + \nub  \mb \nb \aves{B \upara{b}}
%
  - \nun{i} \mi \ni \aves{B \upara{i}}
\nonumber\\[3mm]&&{}
%
  - \nuCX \mi \ni \aves{B \upara{i}}
%
  - \nuOL \mi \ni \aves{B \upara{i}}
\nonumber\\[3mm]&&{}
%
  - \frac{1}{V'} \ddrho V' \aves{|\nabla \rho|} v_{\r{i}\zeta} \mi \ni \aves{B \upara{i}}
%
  + \aves{|\nabla \rho|} v_{\r{i}\zeta} \mi \ni 
  \biggl(\ddrho \frac{\aves{B^2}}{I} \biggl) \frac{I}{\aves{B^2}} \aves{B \upara{i}}
\nonumber\\[3mm]&&{}
%
  + \frac{1}{V'} \ddrho V' \aves{|\nabla \rho|^2} \chis{i\zeta} \mi \ni
  \frac{\aves{B^2}}{I} \ddrho \frac{I}{\aves{B^2}} \aves{B \upara{i}}
\nonumber\\[3mm]&&{}
%
  - \aves{|\nabla \rho|^2} \chis{i\zeta} \mi \ni \biggl( \ddrho
  \frac{\aves{B^2}}{I} \biggr) \ddrho \frac{I}{\aves{B^2}} \aves{B \upara{i}}
%
  - \frac{\aves{B^2}}{I} \frac{1}{V'} \ddrho V' \aves{\Pi_\r{i}^\r{res}}
\nonumber\\[3mm]&&{}
%
  - \frac{1}{V'} \ddrho V' 
    \left[  \aves{|\nabla\rho|} v_{\r{i}\zeta} \mi \ni
          \left(\aves{R^2} - \frac{I^2}{\aves{B^2}}\right) \frac{1}{I}BV_{1\r{i}} \right.
\nonumber\\[3mm]&&{}
    \left.- \aves{|\nabla\rho|^2} \chis{i\zeta} \mi \ni
          \ddrho \left(\aves{R^2} - \frac{I^2}{\aves{B^2}}\right) \frac{1}{I}BV_{1\r{i}}\right]
\\[5mm]
%==============%
% ion toroidal %
%==============%
 \frac{1}{V'} \ddt (V' \mi \ni \aves{R \uzt{i}}) &=&
%
  - \frac{1}{V'} \ddrho V' \left[  \aves{|\nabla\rho|} v_{\r{i}\zeta} \mi \ni \aves{R \uzt{i}}
			       +   (\urho{i} - \ugrho) \mi \ni \aves{R \uzt{i}} \right.
\nonumber\\[3mm]&&{}
                        \left. - \aves{|\nabla\rho|^2} \chis{i\zeta} \mi \ni
			         \ddrho \aves{R \uzt{i}}
                               + \aves{\Pi_\r{i}^\r{res}} \right]
\nonumber\\[3mm]&&{}
%
  + \ell_{11}^\r{ie} \aves{R \uzt{e}}
  + \ell_{11}^\r{ii} \aves{R \uzt{i}}
  + \ell_{11}^\r{iz} \aves{R \uzt{z}}
  + \ell_{11}^\r{ib} \aves{R \uzt{b}}
\nonumber\\[2.5mm]&&{}
%
  - \ell_{12}^\r{ie} \aves{R \qzt{e}}
  - \ell_{12}^\r{ii} \aves{R \qzt{i}}
  - \ell_{12}^\r{iz} \aves{R \qzt{z}}
\nonumber\\[2.5mm]&&{}
%
  + \Zi e                 \ni \aves{R E_\zeta}
%
  + \Zi e \dd{\psi}{\rho} \ni \urho{i}
%
\nonumber\\[2.5mm]&&{}
%
%  turbulent particle flux model
%
%  + \Zi \frac{e^2 \aves{|\nabla\rho|^2} \De}{\Te} \dd{\psi}{\rho} (\ni+\nb)
%    \left[  \dd{\psi}{\rho} \frac{\aves{B^2}}{\bri} \aves{R \uzt{e}}
%          - \dd{\psi}{\rho} \frac{I}{\bri} \aves{B \upara{e}} \right.
%\nonumber\\[2.5mm]&&{}
%    \left.+ \frac{C_\r{T}-1}{e} \ddrho \Te
%	  + \ddrho \Phi
%	  + \frac{\Te}{e} \dd{\psi}{\rho} C_\r{p} \right]
%\nonumber\\[2.5mm]&&{}
%
  + \Zi \frac{e \aves{|\nabla\rho|^2} \De}{\Te} \dd{\psi}{\rho} \frac{\ni+\nb}{\nee}
  \left[  (1-C_\r{T})\Te \ddrho \nee
        + C_\r{T} \ddrho \nee \Te
	+ \nee \Te \dd{\psi}{\rho} C_\r{p}\right]
\nonumber\\[2.5mm]&&{}
%
  - \nuL  \mi \ni \aves{R \uzt{i}}
%
%  + \nub  \mb \frac{I}{\aves{B^2}} \nb \aves{B \upara{b}}
  + \nub  \mb \nb \aves{R \uzt{b}}
%
  - \nun{i} \mi \ni \aves{R \uzt{i}}
%
  - \nuCX \mi \ni \aves{R \uzt{i}}
\nonumber\\[2.5mm]&&{}
%
  - \nuOL \mi \ni \aves{R \uzt{i}}
\\[5mm]
%==========%
% \ni \Ti  %
%==========%
 \frac{3}{2}\frac{1}{V'^{5/3}} \ddt (V'^{5/3} \ni \Ti ) &=& 
%
  - \frac{1}{V'} \ddrho V'
    \left[ - \chis{i}          \aves{|\nabla\rho|^2} \ni \ddrho \Ti
           + V_\r{i}^\r{hp} \aves{|\nabla\rho|}   \ni \Ti
           + \frac{5}{2} \ni \Ti ( \urho{i} - \ugrho ) \right]
\nonumber\\[2.5mm]&&{}
%
  + ( \urho{i} - \ugrho ) \ddrho \ni \Ti
%
  + \uhatth{i} \aves{\vec{B}\cdot\nabla\cdot\tensor{\pi}_\r{i}}
%
  - \frac{3}{2} \ni \frac{ \Ti - \Te }{\tau_\r{ie}}
  - \frac{3}{2} \ni \frac{ \Ti - \Tz }{\tau_\r{iz}}
\nonumber\\[2.5mm]&&{}
%
  - \nuL \Ti (\ni - \nsup{i})
%
  - \frac{3}{2} \nuLT{i} \ni (\Ti - \Tsup{i})
%
\nonumber\\[2.5mm]&&{}
  + \frac{3}{2} \nuion\, \frac{\nee}{\nna + \nnb + \nnc}\, (\nna \Tna +
  \nnb \Tnb + \nnc \Tnc)
\nonumber\\[2.5mm]&&{}
%
  - \frac{3}{2} \nuCX \ni (\Ti - \Tna)
%
  + (1 - \fei) \Pb
%
  + \mb \frac{\Pb}{\Eb} \aves{B \upara{b}} \frac{\aves{B \upara{i}}}{\aves{B^2}}
%
  + \Palp{i}
%
  + \PRF{i}
\\[5mm]
%=====================%
% 2<B \qipara>/(5 pi) %
%=====================%
 \frac{5}{2} \mi \ni \ddt \aves{B \qhatpara{i}} &=&
%
  - \hat{\mu}_2^\r{i} \aves{B \upara{i}} 
  - \hat{\mu}_3^\r{i} \aves{B \qhatpara{i}}
%
  + \hat{\mu}_2^\r{i} {B V_{1\r{i}}} 
  + \hat{\mu}_3^\r{i} {B V_{2\r{i}}}
\nonumber\\[3mm]&&{}
%
  - \ell_{21}^{\r{ie}} \aves{B \upara{e}}
  - \ell_{21}^{\r{ii}} \aves{B \upara{i}}
  - \ell_{21}^{\r{iz}} \aves{B \upara{z}}
  - \ell_{21}^{\r{ib}} \aves{B \upara{b}}
\nonumber\\[3mm]&&{}
%
  + \ell_{22}^{\r{ie}} \aves{B \qhatpara{e}}
  + \ell_{22}^{\r{ii}} \aves{B \qhatpara{i}}
  + \ell_{22}^{\r{iz}} \aves{B \qhatpara{z}}
\nonumber\\[3mm]&&{}
%
  - \nuL \mi \ni \aves{B \qhatpara{i}}
%
  - \nun{i} \mi \ni \aves{B \qhatpara{i}}
\end{eqnarray}
%

%
\begin{eqnarray}
%======%
% \nz  %
%======%
  \frac{1}{V'} \ddt ( V' \nz ) &=&
%
  - \frac{1}{V'} \ddrho [ V' \nz (\urho{z} - \ugrho) ]
%
  + \nee \nnz f_5S_5 - \nee \nz A_6 
\nonumber\\[2.5mm]&&{}
%
  - \nuL (\nz - \nsup{z}) 
%
  + S_\r{LC}
%
  - \nuOL \nz
%
\\[5mm]
%=========%
% radial  %
%=========%
  0 &=&
%
  -           ( \bri ) \ddrho \nz \Tz
%
  - \Zz e \nz ( \bri ) \ddrho \Phi
\nonumber\\[3mm]&&{}
%
  + \Zz e \nz \dd{\psi}{\rho} I          \aves{B\upara{z}}
%
  - \Zz e \nz \dd{\psi}{\rho} \aves{B^2} \aves{R\uzt{z}}
\\[5mm]
%===============%
% <B \upara{z}> %
%===============%
  \mz \nz \ddt \aves{B \upara{z}} &=&
%
  - \hat{\mu}_1^\r{z} \aves{B \upara{z}} 
  - \hat{\mu}_2^\r{z} \aves{B \qhatpara{z}}
%
  + \hat{\mu}_1^\r{z} {B V_{1\r{z}}} 
  + \hat{\mu}_2^\r{z} {B V_{2\r{z}}}
\nonumber\\[3mm]&&{}
%
  + \ell_{11}^{\r{ze}} \aves{B \upara{e}}
  + \ell_{11}^{\r{zi}} \aves{B \upara{i}}
  + \ell_{11}^{\r{zz}} \aves{B \upara{z}}
  + \ell_{11}^{\r{zb}} \aves{B \upara{b}}
\nonumber\\[3mm]&&{}
%
  - \ell_{12}^{\r{ze}} \aves{B \qhatpara{e}}
  - \ell_{12}^{\r{zi}} \aves{B \qhatpara{i}}
  - \ell_{12}^{\r{zz}} \aves{B \qhatpara{z}}
\nonumber\\[3mm]&&{}
%
  + \Zz e \nz \aves{B E_\para}
%
%  + \aves{B S_{\r{m}\para}^\r{i}}
%
%  + \aves{B S_\r{LC\para}}
%
  - \nuL  \mz \nz \aves{B \upara{z}}
%
  - \nun{z} \mz \nz \aves{B \upara{z}}
%
  - \nuOL \mz \nz \aves{B \upara{z}}
\nonumber\\[3mm]&&{}
%
  - \frac{1}{V'} \ddrho V' \aves{|\nabla \rho|} v_{\r{z}\zeta} \mz \nz \aves{B \upara{z}}
%
  + \aves{|\nabla \rho|} v_{\r{z}\zeta} \mz \nz 
  \biggl(\ddrho \frac{\aves{B^2}}{I} \biggl) \frac{I}{\aves{B^2}} \aves{B \upara{z}}
\nonumber\\[3mm]&&{}
%
  + \frac{1}{V'} \ddrho V' \aves{|\nabla \rho|^2} \chis{z\zeta} \mz \nz
  \frac{\aves{B^2}}{I} \ddrho \frac{I}{\aves{B^2}} \aves{B \upara{z}}
\nonumber\\[3mm]&&{}
%
  - \aves{|\nabla \rho|^2} \chis{z\zeta} \mz \nz \biggl( \ddrho
  \frac{\aves{B^2}}{I} \biggr) \ddrho \frac{I}{\aves{B^2}} \aves{B \upara{z}}
%
  - \frac{\aves{B^2}}{I} \frac{1}{V'} \ddrho V' \aves{\Pi_\r{z}^\r{res}}
\nonumber\\[3mm]&&{}
%
  - \frac{1}{V'} \ddrho V' 
    \left[  \aves{|\nabla\rho|} v_{\r{z}\zeta} \mz \nz
          \left(\aves{R^2} - \frac{I^2}{\aves{B^2}}\right) \frac{1}{I}BV_{1\r{z}} \right.
\nonumber\\[3mm]&&{}
    \left.- \aves{|\nabla\rho|^2} \chis{z\zeta} \mz \nz
          \ddrho \left(\aves{R^2} - \frac{I^2}{\aves{B^2}}\right) \frac{1}{I}BV_{1\r{z}}\right]
\\[5mm]
%==============%
% imp toroidal %
%==============%
 \frac{1}{V'} \ddt (V' \mz \nz \aves{R \uzt{z}}) &=&
%
  - \frac{1}{V'} \ddrho V' \left[  \aves{|\nabla\rho|} v_{\r{z}\zeta} \mz \nz \aves{R \uzt{z}}
			       +   (\urho{z} - \ugrho) \mz \nz \aves{R \uzt{z}} \right.
\nonumber\\[3mm]&&{}
                        \left. - \aves{|\nabla\rho|^2} \chis{z\zeta} \mz \nz
			         \ddrho \aves{R \uzt{z}}
                               + \aves{\Pi_\r{z}^\r{res}} \right]
\nonumber\\[3mm]&&{}
%
  + \ell_{11}^\r{ze} \aves{R \uzt{e}}
  + \ell_{11}^\r{zi} \aves{R \uzt{i}}
  + \ell_{11}^\r{zz} \aves{R \uzt{z}}
  + \ell_{11}^\r{zb} \aves{R \uzt{b}}
\nonumber\\[2.5mm]&&{}
%
  - \ell_{12}^\r{ze} \aves{R \qzt{e}}
  - \ell_{12}^\r{zi} \aves{R \qzt{i}}
  - \ell_{12}^\r{zz} \aves{R \qzt{z}}
\nonumber\\[2.5mm]&&{}
%
  + \Zz e                 \nz \aves{R E_\zeta}
%
  + \Zz e \dd{\psi}{\rho} \nz \urho{z}
%
\nonumber\\[2.5mm]&&{}
%
%  turbulent particle flux model
%
%  + \Zz \frac{e \aves{|\nabla\rho|^2} \De}{\Te} \dd{\psi}{\rho} \nz
%    \left[  (1-C_\r{T})\frac{\Te}{\nee}\ddrho \nee
%           + C_\r{T} \frac{1}{\nee} \ddrho \nee\Te
%	  + \Te \dd{\psi}{\rho} C_\r{p} \right]
%\nonumber\\[2.5mm]&&{}
%
  + \Zz \frac{e \aves{|\nabla\rho|^2} \De}{\Te} \dd{\psi}{\rho} \frac{\nz}{\nee}
  \left[  (1-C_\r{T})\Te \ddrho \nee
        + C_\r{T} \ddrho \nee \Te
	+ \nee \Te \dd{\psi}{\rho} C_\r{p}\right]
\nonumber\\[2.5mm]&&{}
%
  - \nuL    \mz \nz \aves{R \uzt{z}}
%
  - \nun{z} \mz \nz \aves{R \uzt{z}}
%
  - \nuOL   \mz \nz \aves{R \uzt{z}}
\\[5mm]
%==========%
% \nz \Tz  %
%==========%
 \frac{3}{2}\frac{1}{V'^{5/3}} \ddt (V'^{5/3} \nz \Tz ) &=& 
%
  - \frac{1}{V'} \ddrho V'
    \left[ - \chis{z}          \aves{|\nabla\rho|^2} \nz \ddrho \Tz
           + V_\r{z}^\r{hp} \aves{|\nabla\rho|}   \nz \Tz
           + \frac{5}{2} \nz \Tz ( \urho{z} - \ugrho ) \right]
\nonumber\\[2.5mm]&&{}
%
  + ( \urho{z} - \ugrho ) \ddrho \nz \Tz
%
  + \uhatth{z} \aves{\vec{B}\cdot\nabla\cdot\tensor{\pi}_\r{z}}
%
  - \frac{3}{2} \nz \frac{ \Tz - \Te }{\tau_\r{ze}}
  - \frac{3}{2} \nz \frac{ \Tz - \Ti }{\tau_\r{zi}}
\nonumber\\[2.5mm]&&{}
%
  - \nuL \Tz (\nz - \nsup{z})
%
  - \frac{3}{2} \nuLT{z} \nz (\Tz - \Tsup{z})
\nonumber\\[2.5mm]&&{}
%
  + \Palp{z}
%
  + \PRF{z}
\\[5mm]
%=====================%
% 2<B \qzpara>/(5 pz) %
%=====================%
 \frac{5}{2} \mz \nz \ddt \aves{B \qhatpara{z}} &=&
%
  - \hat{\mu}_2^\r{z} \aves{B \upara{z}} 
  - \hat{\mu}_3^\r{z} \aves{B \qhatpara{z}}
%
  + \hat{\mu}_2^\r{z} {B V_{1\r{z}}} 
  + \hat{\mu}_3^\r{z} {B V_{2\r{z}}}
\nonumber\\[3mm]&&{}
%
  - \ell_{21}^{\r{ze}} \aves{B \upara{e}}
  - \ell_{21}^{\r{zi}} \aves{B \upara{i}}
  - \ell_{21}^{\r{zz}} \aves{B \upara{z}}
  - \ell_{21}^{\r{zb}} \aves{B \upara{b}}
\nonumber\\[3mm]&&{}
%
  + \ell_{22}^{\r{ze}} \aves{B \qhatpara{e}}
  + \ell_{22}^{\r{zi}} \aves{B \qhatpara{i}}
  + \ell_{22}^{\r{zz}} \aves{B \qhatpara{z}}
\nonumber\\[3mm]&&{}
%
  - \nuL \mz \nz \aves{B \qhatpara{z}}
%
  - \nun{z} \mz \nz \aves{B \qhatpara{z}}
\end{eqnarray}
%

%
\begin{eqnarray}
%======%
% \nb  %
%======%
  \frac{1}{V'} \ddt (V' \nb) &=&
%
  - \frac{1}{V'} \ddrho [ V' \nb (\urho{b} - \ugrho) ]
%
  + \frac{\Pb}{\Eb}
%
  - g \frac{\Pbperp}{\Eb}
%
  - \nub \nb
%
  - \nubL \nb
%
  - \nueff g (\sqrt{\delta} \nb - \nbrp)
%
\nonumber\\[2.5mm]&&{}
  + \frac{1}{V'} \ddrho V'\aves{|\nabla\rho|^2} f_t \Db \ddrho \nb
%
\\[5mm]
%=========%
% radial  %
%=========%
  0 &=&
%
  -           ( \bri ) \ddrho \nb \Tb
%
  - \Zb e \nb ( \bri ) \ddrho \Phi
\nonumber\\[3mm]&&{}
%
  + \Zb e \nb \dd{\psi}{\rho} I          \aves{B\upara{b}}
%
  - \Zb e \nb \dd{\psi}{\rho} \aves{B^2} \aves{R\uzt{b}}
\\[5mm]
%===========%
% \upara{b} %
%===========%
  \mb \ddt \nb \aves{B \upara{b}} &=&
%
  - (\hat{\mu}_1^\r{b} - \ell_{11}^\r{bb}) \aves{B \upara{b}}
%
  + \sum_{a\neq\r{b}} \ell_{11}^{\r{b}a} \aves{B u_{a\para}}
%
  + \hat{\mu}_1^\r{b} {B V_{1\r{b}}}
%
  + \Zb e \nb \aves{B E_\para^A}
\nonumber\\[2.5mm]&&{}
%
  - \nubL \mb \nb \aves{B\upara{b}}
%
  - \nub  \mb \nb \aves{B\upara{b}}
%
  - \nun{b} \mb \nb \aves{B\upara{b}}
%
  - \nuCX \mb \nb \aves{B\upara{b}}
%
  + \aves{B S_{\r{m}\para\r{b}}}
%
%  + \vbph \frac{\Pbpara}{\Eb}
%
\\[5mm]
%===============%
% beam toroidal %
%===============%
  \frac{1}{V'} \ddt (V' \mb \nb \aves{R \uzt{b}}) &=&
%
  - \frac{1}{V'} \dd{}{\rho}[ V' (\urho{b} - \ugrho) \mb \nb \aves{R \uzt{b}} ]
%
%  - \ddV [\aves{|\nabla V|} v_{\r{b}\zeta} \red{\nb \aves{R \uzt{b}}}]
%\nonumber\\[3mm]&&{}
%  + \ddV \aves{|\nabla V|^2} \chis{b\zeta} \ddV \red{\nb \aves{R \uzt{b}}}
%  - \ddV \aves{|\nabla V|^2} \chis{b\zeta} \aves{R \uzt{b}}\ddV \red{\nb} 
%  - \frac{1}{\mb} \ddV \aves{\Pi_\r{b}^\r{res}}
\nonumber\\[2.5mm]&&{}
%
  + \ell_{11}^\r{be} \aves{R \uzt{e}}
  + \ell_{11}^\r{bi} \aves{R \uzt{i}}
  + \ell_{11}^\r{bz} \aves{R \uzt{z}}
%  + \hat{\ell}_{11}^\r{bb} \red{\ni \aves{R \uzt{b}}}
  + \ell_{11}^\r{bb} \nb \aves{R \uzt{b}}
\nonumber\\[2.5mm]&&{}
%
%  - \frac{1}{\mb} \ell_{12}^\r{be} \frac{I}{\aves{B^2}} \red{\aves{B \qhatpara{e}}}
%  - \frac{1}{\mb} \ell_{12}^\r{bi} \frac{I}{\aves{B^2}} \red{\aves{B \qhatpara{i}}}
%
  + \Zb e              \nb \dot{\psi}
%
  + \Zb e \dd{\psi}{\rho} \nb \urho{b}
%
%%
%%  turbulent particle flux model
%%
%  + \Zi \frac{e^2 \aves{|\nabla V|^2} \De}{\mi\Te} \dd{\psi}{V} 
%%
%    \left[  \dd{\psi}{V} \frac{\aves{B^2}}{\bri} \red{\nee \aves{R \uzt{e}}}
%          - \dd{\psi}{V} \frac{I}{\bri} \nee \red{\aves{B \upara{e}}} \right.
%\nonumber\\[2.5mm]&&{}
%%
%    \left.+ \frac{C_\r{T}-1}{e} \left(\ddV \red{\nee \Te} - \Te \ddV \red{\nee} \right)
%	  + \nee \ddV \red{\Phi}
%	  + \frac{\Te}{e} \dd{\psi}{V} C_\r{p} \red{\nee} \right]
%\nonumber\\[2.5mm]&&{}
%
  - \nubL \mb \nb \aves{R \uzt{b}}
%
  - \nub  \mb \nb \aves{B \uzt{b}}
\nonumber\\[2.5mm]&&{}
%
  - \nun{i} \mb \nb \aves{R \uzt{b}}
%
  - \nuCX \mb \nb \aves{R \uzt{b}}
%
%  - \nuOL \mb \nb \aves{R \uzt{b}}
%
  + \frac{I}{\aves{B^2}} \aves{B S_{\r{m}\para\r{b}}}
%
\\[5mm]
%=======%
% \nna  %
%=======%
  \frac{1}{V'} \ddt (V' \nna) &=& 
%
    \frac{1}{V'} \ddrho V' \aves{|\nabla\rho|^2} D_{01} \ddrho \nna
%
  - \nuion\, \frac{\nee}{\nna + \nnb + \nnc}\, \nna
%
\nonumber\\[2.5mm]&&{}
  - \nuCX\, \frac{\ni}{\nna + \nnb}\, \nna
%
  + \gamma_0 \nuL \, (\ni - \nsup{i})
%
  + \frac{1}{V'} \ddrho V' \aves{\vec{\Gamma}^\r{puff}_\r{i}\cdot\nabla\rho}\biggr|_\r{b}
\\[5mm]
%=======%
% \nnb  %
%=======%
  \frac{1}{V'} \ddt (V' \nnb) &=& 
%
    \frac{1}{V'} \ddrho V' \aves{|\nabla\rho|^2} D_{02} \ddrho \nnb
%
  - \nuion\, \frac{\nee}{\nna + \nnb + \nnc}\, \nnb
%
%\nonumber\\[2.5mm]&&{}
  + \nuCX\, \frac{\ni}{\nna + \nnb}\, \nna
%
%  + \fCX \frac{\Pb}{\Eb}
\\[5mm]
%=======%
% \nnc  %
%=======%
  \frac{1}{V'} \ddt (V' \nnc) &=& 
%
    \frac{1}{V'} \ddrho V' \aves{|\nabla\rho|^2} D_{03} \ddrho \nnc
%
  - \nuion\, \frac{\nee}{\nna + \nnb + \nnc}\, \nnc
%
%\nonumber\\[2.5mm]&&{}
  + \fCX \frac{\Pb}{\Eb}
\\[5mm]
%=======%
% \nnz  %
%=======%
  \frac{1}{V'} \ddt (V' \nnz) &=& 
%
    \frac{1}{V'} \ddrho V' \aves{|\nabla\rho|^2} D_{n\r{z}} \ddrho \nnz
%
  + \frac{1}{V'} \ddrho V' \aves{|\nabla\rho|^2} \Delta D_{\r{C0}} \ddrho n_\r{C0}
%
  - \nee \nnz f_5S_5 + \nee \nz A_6 
\nonumber\\[2.5mm]&&{}
%
  + \gamma_\r{0z} \nuL \, (\nz - \nsup{z})
%
  + \gamma_\r{i,z} \nuL \, (\ni - \nsup{i})
%
  - (1-\gamma_\r{0z}) \nuL \, (\nnz - n_\r{C0})
\nonumber\\[2.5mm]&&{}
%
  + \frac{1}{V'} \ddrho V' \aves{\vec{\Gamma}^\r{puff}_\r{z}\cdot\nabla\rho}\biggr|_\r{b}
\\[5mm]
%========%
% \nbrp  %
%========%
  \frac{1}{V'} \ddt (V' g \nbrp) &=&
%
    g \frac{\Pbperp}{\Eb}
%
  + g \nueff (\sqrt{\delta} \nb - \nbrp)
%
  - g \nub \nbrp
%
\nonumber\\[2.5mm]&&{}
  - \frac{1}{V'} \ddrho V' g \left( \aves{|\nabla\rho|} \ubrp \nbrp - \aves{|\nabla\rho|^2} \Dbrp \ddrho \nbrp \right)
%
\end{eqnarray}

\bigskip

\section{Additional important relationship}
\setcounter{equation}{0}

\begin{eqnarray}
%============================================%
% superstaged impurity w/o fully striped one %
%============================================%
  \nnz &=&
%
 \sum_{i=0}^{\r{max}(Z)-1} n_i
\\[5mm]
%=======%
% B_pol %
%=======%
  B_p &=&
%
  \frac{|\nabla\rho|}{R} \dd{\psi}{\rho},
  \quad
%===========%
% <B_pol^2> %
%===========%
  \aves{B_p^2} =
%
  \ave{\frac{|\nabla V|^2}{R^2}} \left(\dd{\psi}{V}\right)^2
\\[5mm]
%=======%
% B_tor %
%=======%
  B_t &=&
%
  \frac{I}{R}
%
  = \frac{1}{R} \frac{4\pi^2}{\aves{R^{-2}}} \dd{\psit}{V},
  \quad
%===========%
% <B_tor^2> %
%===========%
  \aves{B_t^2} =
%
  \frac{16\pi^4}{\aves{R^{-2}}} \left(\dd{\psit}{V}\right)^2
\\[5mm]
%=======%
% <B^2> %
%=======%
  \aves{B^2} &=&
%
  \aves{B_t^2} + \aves{B_p^2}
%
  = \frac{16\pi^4}{\aves{R^{-2}}}\left(\dd{\psit}{V}\right)^2
  + \ave{\frac{|\nabla V|^2}{R^2}}\left(\dd{\psi}{V}\right)^2
\\[5mm]
%===========%
% <B^theta> %
%===========%
  \aves{B^\theta} &=&
%
  \frac{4\pi^2}{V'} \dd{\psi}{\rho}
%
  = 4\pi^2 \dd{\psi}{V}
%
  = \frac{I\aves{R^{-2}}}{q}
%
  = 2\pi\left(\oint \frac{dl_p}{B_p}\right)^{-1}
\\[5mm]
%=======%
% q hat %
%=======%
  \hat{q}^2 &\equiv& \frac{I^2}{2\aves{B_p^2}}\left(\ave{\frac{1}{R^2}}-\frac{1}{\aves{R^2}}\right)
\\[5mm]
%================%
% <B^2><R^2>-I^2 %
%================%
  \bri &=&
%
  \aves{R^2} \ave{\frac{|\nabla V|^2}{R^2}} \left(\dd{\psi}{V}\right)^2 (1+2\hat{q}^2)
\\[5mm]
%===================%
% 2<B \qpara>/(5 p) %
%===================%
 \aves{B\hat{q}_{s\para}} &\equiv&
%
 \frac{2\aves{Bq_{s\para}}}{5p_s}
\\[5mm]
%=============%
% hat u_theta %
%=============%
  \hat{u}_{s\theta} &=&
%
  \frac{\aves{B\uspara} - BV_{1s}}{ \aves{B^2} }
\\[5mm]
%=====================%
% <uzt/R> and <R uzt> %
%=====================%
  \ave{\frac{\uszt}{R}} &=&
%
  \frac{1}{1+2\hat{q}^2}\frac{1}{\aves{R^2}}\aves{R\uszt}+\frac{2\hat{q}^2}{1+2\hat{q}^2}\frac{1}{I}\aves{B\uspara}
\\[2.5mm]
%
%  &=& \frac{\aves{R\uszt}}{\aves{R^2}}
%
%  + \frac{2\hat{q}^2}{1+2\hat{q}^2} \left( \frac{1}{I}\aves{B\uspara} - \frac{1}{\aves{R^2}}\aves{R\uszt}\right)
%\nonumber\\[2.5mm]
%
  &=& \frac{\aves{R\uszt}}{\aves{R^2}}
%
  + \frac{I}{\aves{B^2}} \left(  \ave{\frac{1}{R^2}}
			      - \frac{1}{\aves{R^2}} \right)
    (\aves{B\uspara} - BV_{1s})
\nonumber\\[2.5mm]
%
  &=& \frac{\aves{R\uszt}}{\aves{R^2}}
%
  + I \left(  \ave{\frac{1}{R^2}}
			       - \frac{1}{\aves{R^2}} \right) \hat{u}_{s\theta}
\nonumber\\[5mm]
%=========%
% <R qzt> %
%=========%
  \ave{R \qzt{s}} &=&
%
    \frac{1}{\aves{B^2}} \left( \bri \right) \frac{1}{I} B V_{2s}
%
  + \frac{I}{\aves{B^2}} \aves{B\qhatpara{s}}
\\[5mm]
%=======%
% B V_1 %
%=======%
  B V_{1s} &=&
%
 - I \left( \frac{1}{e_s n_s} \dd{n_s T_s}{\psi} + \dd{\Phi}{\psi}\right),
\quad
%\\[5mm]
%=======%
% B V_2 %
%=======%
  B V_{2s} =
%
 - \frac{I}{e_s} \dd{T_s}{\psi}
\\[5mm]
%====================%
% u.grad rho = u^rho %
%====================%
  u^\rho &\equiv&
%
  \aves{ \vec{u} \cdot \nabla\rho },
%
  \quad u^V = \aves{ \vec{u} \cdot \nabla V }
\\[5mm]
%===============%
% grid velocity %
%===============%
  \ugrho &=&
%
  - \frac{4\pi^2}{IV'\aves{R^{-2}}} \psitd,
%
  \quad \ugV = - \frac{4\pi^2}{I\aves{R^{-2}}} \psitd = - \frac{V'}{2\rho\psi_{ta}}\psitd
\\[5mm]
  \frac{1}{V'}\dd{V'}{t} &=&
%
  \dd{}{V}\aves{\vec{u}_g\cdot\nabla V} = \dd{u_g^V}{V}
\\[5mm]
%======%
% \Er  %
%======%
  \aves{\Er} &=&
%
  - \aves{|\nabla\rho|^2} \dd{r}{\rho} \ddrho \Phi
\\[5mm]
%============%
% <B E_para> %
%============%
  \aves{B E_\para} &=&
%
  I \aves{R^{-2}} \left( - \frac{\psitd}{q} + \psid \right)
\\[5mm]
%=========%
% R E_tor %
%=========%
  R E_\zeta &=&
%
  \psid
%
  = \frac{V_\r{loop}}{2\pi}
\\[5mm]
%=================%
% viscous heating %
%=================%
  \aves{\vec{B}\cdot\nabla\cdot\tensor{\pi}_\r{i}} &=&
%
  \hat{\mu}_1^\r{i} ( \aves{Bu_{\r{i}\para}} - BV_{1\r{i}}) 
%
  + \hat{\mu}_2^\r{i} \left( \frac{2\aves{Bq_{\r{i}\para}}}{5p_\r{i}} -
		       BV_{2\r{i}}\right)
\\[5mm]
%======================%
% friction coefficient %
%======================%
  \hat{\ell}_{ij}^{ab} &\equiv& \frac{\ell_{ij}^{ab}}{m_an_a},
  \quad
%===========================%
% beam friction coefficient %
%===========================%
  \hat{\ell}_{ij}^{\r{b}b} \equiv \frac{\ell_{ij}^{\r{b}b}}{m_\r{e}n_\r{e}}
\\[5mm]
%==========%
% gas puff %
%==========%
  \Gamma^\r{puff} &\simeq& \frac{\aves{\vec{\Gamma}^\r{puff}\cdot\nabla\rho}}{\aves{|\nabla\rho|}}
\end{eqnarray}

\bigskip

\section{Equations in the coordinates suitable for implementation}
\setcounter{equation}{0}

\begin{eqnarray}
%=======%
% \phi  %
%=======%
  0 &=&
%
    \ddV \left[ \aves{|\nabla V|^2} \ddV{\red{\Phi}} \right]
%
  + \frac{1}{\epsilon_0}
%
    \left(-e \red{\nee} + \Zi e \red{\ni} + \Zz e \red{\nz} + \Zb e \red{\nb} + \Zb e g \red{\nbrp} \right)
\\[5mm]
%=========%
% \psitd  %
%=========%
%  \frac{1}{c^2} \ave{\frac{|\nabla V|^2}{R^2}}^{-1} \ddt \red{\psitd} &=&
%%
%   \ddV \left[ \frac{1}{\aves{R^{-2}}}
%				         \ddV \red{\psit} \right]
%%
%  + \frac{\mu_0}{\aves{B^\theta}} \ave{\frac{|\nabla V|^2}{R^2}}^{-1}
% \nonumber\\[3mm]&&{}
%%
%    \times \left[-     e \left( \nee \red{\aves{B\upara{e}}} - I\red{\nee \ave{\frac{\uzt{e}}{R}}} \right)
%          + \Zi e \left( \ni \red{\aves{B\upara{i}}} - I\red{\ni  \ave{\frac{\uzt{i}}{R}}} \right)
%    \right.\nonumber\\[3mm]&&{}
%%    \left.+ \Zb e  \left( \red{\nb \aves{B\upara{b}}} - I\red{\nb \ave{\frac{\uzt{b}}{R}}} \right) \right]
%    \left.+ \Zb e  \frac{1}{\aves{B^2}} \ave{\frac{|\nabla V|^2}{R^2}}
%     \left(\dd{\psi}{V}\right)^2 \nb \red{\aves{B\upara{b}}} \right]
%
  \epsilon_0 \ddt \red{\psitd} &=&
%
     \ddV \left[ \ave{\frac{|\nabla V|^2}{R^2}} \frac{1}{\aves{R^{-2}}} \ddV \red{\frac{\psit}{\mu_0}} \right]
   - \frac{1}{\aves{R^{-2}}} \ddV \left[ \ave{\frac{|\nabla V|^2}{R^2}}
				  \right] \ddV  \red{\frac{\psit}{\mu_0}}
 \nonumber\\[3mm]&&{}
%
  + \frac{1}{\aves{B^\theta}}
%
    \left[-     e \left( \nee \red{\aves{B\upara{e}}} - I\red{\nee \ave{\frac{\uzt{e}}{R}}} \right)
          + \Zi e \left( \ni  \red{\aves{B\upara{i}}} - I\red{\ni  \ave{\frac{\uzt{i}}{R}}} \right)
    \right.\nonumber\\[3mm]&&{}
    \left.+ \Zz e \left( \nz  \red{\aves{B\upara{z}}} - I\red{\nz  \ave{\frac{\uzt{z}}{R}}} \right)
          + \Zb e  \left( \red{\nb \aves{B\upara{b}}} - I\red{\nb  \ave{\frac{\uzt{b}}{R}}} \right) \right]
%    \left.+ \Zb e  \frac{1}{\aves{B^2}} \ave{\frac{|\nabla V|^2}{R^2}}
%     \left(\dd{\psi}{V}\right)^2 \red{\nb \aves{B\upara{b}}} \right]
\\[5mm]
%========%
% \psid  %
%========%
  \frac{1}{c^2} \ave{\frac{1}{R^2}} \ddt \red{\psid} &=&
%
  \ddV \left[ \ave{\frac{|\nabla V|^2}{R^2}}
				    \ddV \red{\psi} \right] \nonumber\\[3mm]&&{}
%
  - \mu_0 \left(-     e \red{\nee \ave{\frac{\uzt{e}}{R}}}
                + \Zi e \red{\ni  \ave{\frac{\uzt{i}}{R}}}
                + \Zz e \red{\nz  \ave{\frac{\uzt{z}}{R}}}
	        + \Zb e \red{\nb  \ave{\frac{\uzt{b}}{R}}} \right)
%	        + \Zb e \ave{\frac{1}{R^2}} \frac{I}{\aves{B^2}}
%	        \red{\nb \aves{B \upara{b}} }\right)
%
%  + \mu_0 j_\r{virtual}
\\[5mm]
%=======%
% \psi  %
%=======%
  \ddt \red{\psi} &=&
%
  \red{\psid}
\\[5mm]
%=========%
% \psitd  %
%=========%
  \ddt \red{\frac{\psit}{\mu_0}} &=&
%
  \frac{\red{\psitd}}{\mu_0}
\end{eqnarray}
%

%
\begin{eqnarray}
%=======%
% \nee  %
%=======%
  \ddt \red{\nee}  &=&
%
  - \ddV \red{\nee\uV{e}}
%
  + \ugV \ddV \red{\nee}
%
  + \nuion\, \frac{\nee}{\nna + \nnb + \nnc}\, (\red{\nna} + \red{\nnb}
  + \red{\nnc})
\nonumber\\[3mm]&&{}
%
  + \Zz \nee \red{\nnz} f_5S_5 - \Zz \nee \red{\nz} A_6
%
  - \nuL (\red{\nee} - \nsup{e})
%
  + (1 - \fCX) \frac{\Pb}{\Eb}
\\[5mm]
%=========%
% radial  %
%=========%
  0 &=&
%
  - \frac{e}{\me} \nee \left( \bri \right) \frac{1}{I} \red{BV_{1\r{e}}}
%
  - \frac{e}{\me} I \nee \red{\aves{B\upara{e}}}
%
  + \frac{e}{\me} \aves{B^2} \red{\nee \aves{R\uzt{e}}}
\\[5mm]
%===============%
% <B \upara{e}> %
%===============%
  \nee \ddt \red{\aves{B \upara{e}}} &=&
%
  - \frac{1}{\me} \hat{\mu}_1^\r{e} \red{\aves{B \upara{e}}}
  - \frac{1}{\me} \hat{\mu}_2^\r{e} \red{\aves{B \qhatpara{e}}}
%
  + \frac{1}{\me} \hat{\mu}_1^\r{e} \red{BV_{1\r{e}}}
\nonumber\\[3mm]&&{}
%
%  + \frac{1}{\me} \hat{\mu}_2^\r{e} {B V_{2\r{e}}}
  + \frac{1}{e\me} \hat{\mu}_2^\r{e} I \dd{V}{\psi} \frac{1}{\nee} \ddV
  \red{\nee \Te}
  - \frac{1}{e\me} \hat{\mu}_2^\r{e} I \dd{V}{\psi} \frac{\Te}{\nee} \ddV
  \red{\nee}
\nonumber\\[3mm]&&{}
%
  + \hat{\ell}_{11}^{\r{ee}} \nee \red{\aves{B \upara{e}}}
  + \hat{\ell}_{11}^{\r{ei}} \nee \red{\aves{B \upara{i}}}
  + \hat{\ell}_{11}^{\r{ez}} \nee \red{\aves{B \upara{z}}}
  + \hat{\ell}_{11}^{\r{eb}} \frac{\nee}{\nb} \red{\nb \aves{B \upara{b}}}
\nonumber\\[3mm]&&{}
%
  - \hat{\ell}_{12}^{\r{ee}} \nee \red{\aves{B \qhatpara{e}}}
  - \hat{\ell}_{12}^{\r{ei}} \nee \red{\aves{B \qhatpara{i}}}
  - \hat{\ell}_{12}^{\r{ez}} \nee \red{\aves{B \qhatpara{z}}}
\nonumber\\[3mm]&&{}
%
  - e \frac{\nee}{\me} I \aves{R^{-2}} \left( - \frac{\red{\psitd}}{q} + \red{\psid} \right)
%
%  + \frac{1}{\me} \aves{B S_{\r{m}\para}^\r{e}}
%
  - \nuL \nee \red{\aves{B \upara{e}}}
%
  - \nun{e} \nee \red{\aves{B \upara{e}}}
\nonumber\\[3mm]&&{}
%
  - \ddV \aves{|\nabla \rho|} v_{\r{e}\zeta} \nee \red{\aves{B \upara{e}}}
%
  + \aves{|\nabla V|} v_{\r{e}\zeta} \nee 
  \biggl(\ddV \frac{\aves{B^2}}{I} \biggl) \frac{I}{\aves{B^2}} \red{\aves{B \upara{e}}}
\nonumber\\[3mm]&&{}
%
  + \ddV \aves{|\nabla V|^2} \chis{e\zeta} \nee
  \frac{\aves{B^2}}{I} \ddV \frac{I}{\aves{B^2}} \red{\aves{B \upara{e}}}
\nonumber\\[3mm]&&{}
%
  - \aves{|\nabla V|^2} \chis{e\zeta} \nee \biggl( \ddV
  \frac{\aves{B^2}}{I} \biggr) \ddV \frac{I}{\aves{B^2}} \red{\aves{B \upara{e}}}
%
  - \frac{1}{\me} \frac{\aves{B^2}}{I} \ddV \aves{\Pi_\r{e}^\r{res}}
\nonumber\\[3mm]&&{}
%
  - \ddV
    \left[  \aves{|\nabla V|} v_{\r{e}\zeta} \nee
            \frac{\aves{R^2}}{\aves{B^2}}\ave{\frac{|\nabla V|^2}{R^2}}\dd{\psi}{V}(1+2\hat{q}^2)
	    \frac{1}{I} \dd{\psi}{V} BV_{1\r{e}} \right.
\nonumber\\[3mm]&&{}
    \left.- \aves{|\nabla V|^2} \chis{e\zeta} \nee
          \ddV \frac{\aves{R^2}}{\aves{B^2}}\ave{\frac{|\nabla V|^2}{R^2}}\dd{\psi}{V}(1+2\hat{q}^2)
	  \frac{1}{I} \dd{\psi}{V} BV_{1\r{e}}\right]
\\[5mm]
%================%
% elec. toroidal %
%================%
 \ddt \red{\nee \aves{R \uzt{e}}} &=&
%
  - (\uV{e} - \ugV) \ddV \red{\nee \aves{R \uzt{e}}}
%
  - \dd{\uV{e}}{V} \red{\nee \aves{R \uzt{e}}}
%
  - \ddV [\aves{|\nabla V|} v_{\r{e}\zeta} \red{\nee \aves{R \uzt{e}}}]
\nonumber\\[3mm]&&{}
  + \ddV \aves{|\nabla V|^2} \chis{e\zeta} \ddV \red{\nee \aves{R \uzt{e}}}
  - \ddV \aves{|\nabla V|^2} \chis{e\zeta} \aves{R \uzt{e}} \ddV \red{\nee}
  - \frac{1}{\me} \ddV \aves{\Pi_\r{e}^\r{res}}
\nonumber\\[2.5mm]&&{}
%
  + \hat{\ell}_{11}^\r{ee}                  \red{\nee \aves{R \uzt{e}}}
  + \hat{\ell}_{11}^\r{ei} \frac{\nee}{\ni} \red{\ni  \aves{R \uzt{i}}}
  + \hat{\ell}_{11}^\r{ez} \frac{\nee}{\nz} \red{\nz  \aves{R \uzt{z}}}
  + \hat{\ell}_{11}^\r{eb} \frac{\nee}{\nb} \red{\nb  \aves{R \uzt{b}}}
%  + \hat{\ell}_{11}^\r{eb} \frac{I}{\aves{B^2}} \frac{\nee}{\nb}
%  \red{\nb \aves{B \upara{b}}}
\nonumber\\[2.5mm]&&{}
%
%  - \hat{\ell}_{12}^\r{ee} \frac{I}{\aves{B^2}} \nee \red{\aves{B \qhatpara{e}}}
%  - \hat{\ell}_{12}^\r{ei} \frac{I}{\aves{B^2}} \nee \red{\aves{B \qhatpara{i}}}
%  - \hat{\ell}_{12}^\r{ez} \frac{I}{\aves{B^2}} \nee \red{\aves{B \qhatpara{z}}}
  - \hat{\ell}_{12}^\r{ee} \nee \frac{1}{e} \frac{1}{\aves{B^2}} \left( \bri \right)
    \dd{V}{\psi} \ddV \Te
  - \hat{\ell}_{12}^\r{ee} \frac{I}{\aves{B^2}} \red{\aves{B \qhatpara{e}}}
\nonumber\\[2.5mm]&&{}
  + \hat{\ell}_{12}^\r{ei} \nee \frac{1}{\Zi e} \frac{1}{\aves{B^2}} \left( \bri \right)
    \dd{V}{\psi} \ddV \Ti
  - \hat{\ell}_{12}^\r{ei} \frac{I}{\aves{B^2}} \red{\aves{B \qhatpara{i}}}
\nonumber\\[2.5mm]&&{}
  + \hat{\ell}_{12}^\r{ez} \nee \frac{1}{\Zz e} \frac{1}{\aves{B^2}} \left( \bri \right)
    \dd{V}{\psi} \ddV \Tz
  - \hat{\ell}_{12}^\r{ez} \nee \frac{I}{\aves{B^2}} \red{\aves{B \qhatpara{z}}}
\nonumber\\[2.5mm]&&{}
%
  - e \frac{\nee}{\me} \red{\psid}
%
  - e \frac{1}{\me} \dd{\psi}{V} \red{\nee \uV{e}}
%
\nonumber\\[2.5mm]&&{}
%p
%%%  turbulent particle flux model
%
  - \frac{e^2 \aves{|\nabla V|^2} \De}{\me\Te} \dd{\psi}{V} 
%%%
%    \left[  \dd{\psi}{V} \frac{\aves{B^2}}{\bri} \red{\nee \aves{R \uzt{e}}}
%%
%          - \dd{\psi}{V} \frac{I}{\bri} \nee \red{\aves{B \upara{e}}} \right.
%\nonumber\\[2.5mm]&&{}
%%
%    \left.+ \frac{C_\r{T}-1}{e} \left( \ddV \red{\nee \Te} - \Te \ddV \red{\nee} \right)
%%
%	  + \nee \ddV \red{\Phi}
%
%    	  + \frac{\Te}{e} \dd{\psi}{V} C_\r{p} \red{\nee} \right]
%%%
  \left[ \nee \dd{\psi}{V} \frac{1}{I} \red{B V_{1\r{e}}}
       + \nee \ddV \red{\Phi}
       + \frac{C_\r{T}-1}{e} \left( \ddV \red{\nee \Te} - \Te \ddV
				    \red{\nee} \right)
       + \frac{\Te}{e} \dd{\psi}{V} C_\r{p} \red{\nee} \right]
\nonumber\\[2.5mm]&&{}
%%%
  - \nuL  \red{\nee \aves{R \uzt{e}}}
%
  - \nun{e} \red{\nee \aves{R \uzt{e}}}
\\[5mm]
%===========%
% \nee \Te  %
%===========%
 \frac{3}{2} \ddt \red{\nee \Te} &=& 
%
  - \frac{5}{2} ( \uV{e} - \ugV ) \ddV \red{\nee \Te}
%
  - \frac{5}{2} \dd{\uV{e}}{V} \red{\nee \Te}
%
  - \ddV [V_\r{e}^\r{hp} \aves{|\nabla V|}   \red{\nee \Te}]
\nonumber\\[2.5mm]&&{}
%
  + \ddV \chis{e} \aves{|\nabla V|^2}\ddV \red{\nee \Te}
  - \ddV \chis{e} \aves{|\nabla V|^2}\Te \ddV \red{\nee}
\nonumber\\[2.5mm]&&{}
%
  - \ugV   \ddV \red{\nee \Te}
  - \uV{i} \ddV \red{\ni  \Ti}
  - \uV{z} \ddV \red{\nz  \Tz}
%
  - \uhatth{i} \aves{\vec{B}\cdot\nabla\cdot\tensor{\pi}_\r{i}}
  - \uhatth{z} \aves{\vec{B}\cdot\nabla\cdot\tensor{\pi}_\r{z}}
\nonumber\\[2.5mm]&&{}
%
  - \frac{3}{2} \nu_\r{Tei} \left(\red{\nee \Te} - \frac{\nee}{\ni} \red{\ni \Ti}\right)
  - \frac{3}{2} \nu_\r{Tez} \left(\red{\nee \Te} - \frac{\nee}{\nz} \red{\nz \Tz}\right)
\nonumber\\[2.5mm]&&{}
%
%  - \dd{V}{\psi} \ddV ( \red{\nee \Te} + \red{\ni \Ti} ) \psit
%
%  - \frac{1}{\mu_0} \dd{V}{\psi} \dd{I}{V} I \aves{R^{-2}}
%
%    \left( - \frac{\red{\psitd}}{q} + \red{\psid} \right)
  - 4\pi^2 \dd{\psi}{V} \dot{\psi}_t 
  (-    e \nee \uhatth{e} + \Zi e \ni  \uhatth{i} + \Zz e \nz  \uhatth{z})
%  - \frac{1}{\aves{B^2}}\dot{\psi}_t (-e \nee \red{\aves{Bu_{\r{e}\para}}} +
%  Z_\r{i} e \ni \red{\aves{Bu_{\r{i}\para}}})
%
%  +\frac{1}{\aves{B^2}}\dot{\psi}_t (-e \nee BV_{1\r{e}} + Z_\r{i} e \ni
%  BV_{1\r{i}})
\nonumber\\[2.5mm]&&{}
%
  +\dot{\psi}\left(-     e \red{\nee \ave{\frac{\uzt{e}}{R}}}
	           + \Zi e \red{\ni  \ave{\frac{\uzt{i}}{R}}}
	           + \Zz e \red{\nz  \ave{\frac{\uzt{z}}{R}}}\right)
\nonumber\\[2.5mm]&&{}
%
  - \nuL (\red{\nee\Te} - \nsup{e}\Te)
%
  - \frac{3}{2} \nuLT{e} (\red{\nee\Te} - \nee\Tsup{e})
%
  - \EH \nuion\, \frac{\nee}{\nna + \nnb + \nnc}\, (\red{\nna} + \red{\nnb} + \red{\nnc})
\nonumber\\[2.5mm]&&{}
%
  + \fei \Pb
%
%  + \mb \frac{1}{\nb} \frac{\Pb}{\Eb} \frac{\aves{B \upara{e}}}{\aves{B^2}} \red{\nb \aves{B \upara{b}}}
%
  + \Palp{e}
%
  + \PRF{e}
%
  - \PBr
\\[5mm]
%=====================%
% 2<B \qepara>/(5 pe) %
%=====================%
 \frac{5}{2} \nee \ddt \red{\aves{B \qhatpara{e}}} &=&
%
  - \frac{1}{\me} \hat{\mu}_2^\r{e} \red{\aves{B \upara{e}}}
  - \frac{1}{\me} \hat{\mu}_3^\r{e} \red{\aves{B \qhatpara{e}}}
%
  + \frac{1}{\me} \hat{\mu}_2^\r{e} \red{B V_{1\r{e}}}
\nonumber\\[3mm]&&{}
%
%  + \frac{1}{\me} \hat{\mu}_3^\r{e} {B V_{2\r{e}}}
  + \frac{1}{e\me} \hat{\mu}_3^\r{e} I \dd{V}{\psi} \frac{1}{\nee} \ddV
  \red{\nee \Te}
  - \frac{1}{e\me} \hat{\mu}_3^\r{e} I \dd{V}{\psi} \frac{\Te}{\nee} \ddV
  \red{\nee}
\nonumber\\[3mm]&&{}
%
  - \hat{\ell}_{21}^{\r{ee}} \nee \red{\aves{B \upara{e}}}
  - \hat{\ell}_{21}^{\r{ei}} \nee \red{\aves{B \upara{i}}}
  - \hat{\ell}_{21}^{\r{ez}} \nee \red{\aves{B \upara{z}}}
  - \hat{\ell}_{21}^{\r{eb}} \frac{\nee}{\nb} \red{\nb \aves{B \upara{b}}}
\nonumber\\[3mm]&&{}
%
  + \hat{\ell}_{22}^{\r{ee}} \nee \red{\aves{B \qhatpara{e}}}
  + \hat{\ell}_{22}^{\r{ei}} \nee \red{\aves{B \qhatpara{i}}}
  + \hat{\ell}_{22}^{\r{ez}} \nee \red{\aves{B \qhatpara{z}}}
\nonumber\\[3mm]&&{}
%
  - \nuL \nee \red{\aves{B \qhatpara{e}}}
%
  - \nun{e} \nee \red{\aves{B \qhatpara{e}}}
\\[5mm]
%=======================%
% <uezt/R> and <R uezt> %
%=======================%
    0 &=&
%
  - \red{\nee \ave{\frac{\uzt{e}}{R}}}
%
  + \frac{1}{\aves{R^2}} \red{\nee \aves{R\uzt{e}}}
%
  - \frac{2\hat{q}^2}{1+2\hat{q}^2} \frac{1}{\aves{R^2}}\red{\nee\aves{R\uzt{e}}}
%
  + \frac{2\hat{q}^2}{1+2\hat{q}^2} \frac{1}{I}\nee \red{\aves{B\upara{e}}}
\\[5mm]
%=======%
% BV_1e %
%=======%
    0 &=&
%
  - \nee \red{BV_{1\r{e}}}
%
  + \frac{1}{e} \dd{V}{\psi} I \ddV \red{\nee\Te}
%
  - \nee \dd{V}{\psi} I \ddV \red{\Phi}
\end{eqnarray}
%

%\clearpage

%
\begin{eqnarray}
%======%
% \ni  %
%======%
  \ddt \red{\ni} &=&
%
  - \ddV \red{\ni\uV{i}}
%
  + \ugV \ddV \red{\ni}
%
  + \nuion\, \frac{\nee}{\nna + \nnb + \nnc}\, (\red{\nna} +
  \red{\nnb} + \red{\nnc})
%
  - \nuL (\red{\ni} - \nsup{i}) 
\nonumber\\[2.5mm]&&{}
%
  + \nub \red{\nb}
%
  + \nub g \red{\nbrp}
%
  - \fCX \frac{\Pb}{\Eb}
%
  + \nubL \red{\nb}
%
  + S_\r{LC}
%
  - \nuOL \red{\ni}
%
\\[5mm]
%=========%
% radial  %
%=========%
  0 &=&
%
    \Zi \frac{e}{\mi} \ni \left( \bri \right) \frac{1}{I} \red{B V_{1\r{i}}}
%
  + \Zi \frac{e}{\mi} I \ni \red{\aves{B\upara{i}}}
%
  - \Zi \frac{e}{\mi} \aves{B^2} \red{\ni \aves{R\uzt{i}}}
\\[5mm]
%===============%
% <B \upara{i}> %
%===============%
  \ni \ddt \red{\aves{B \upara{i}}} &=&
%
  - \frac{1}{\mi} \hat{\mu}_1^\r{i} \red{\aves{B \upara{i}}}
  - \frac{1}{\mi} \hat{\mu}_2^\r{i} \red{\aves{B \qhatpara{i}}}
%
  + \frac{1}{\mi} \hat{\mu}_1^\r{i} \red{B V_{1\r{i}}} 
\nonumber\\[3mm]&&{}
%
%  + \frac{1}{\mi} \hat{\mu}_2^\r{i} {B V_{2\r{i}}}
  - \frac{1}{\Zi e\mi} \hat{\mu}_2^\r{i} I \dd{V}{\psi} \frac{1}{\ni}
  \ddV \red{\ni\Ti}
  + \frac{1}{\Zi e\mi} \hat{\mu}_2^\r{i} I \dd{V}{\psi} \frac{\Ti}{\ni}
  \ddV \red{\ni}
\nonumber\\[3mm]&&{}
%
  + \hat{\ell}_{11}^{\r{ie}} \ni \red{\aves{B \upara{e}}}
  + \hat{\ell}_{11}^{\r{ii}} \ni \red{\aves{B \upara{i}}}
  + \hat{\ell}_{11}^{\r{iz}} \ni \red{\aves{B \upara{z}}}
  + \hat{\ell}_{11}^{\r{ib}} \frac{\ni}{\nb} \red{\nb \aves{B \upara{b}}}
\nonumber\\[3mm]&&{}
%
  - \hat{\ell}_{12}^{\r{ie}} \ni \red{\aves{B \qhatpara{e}}}
  - \hat{\ell}_{12}^{\r{ii}} \ni \red{\aves{B \qhatpara{i}}}
  - \hat{\ell}_{12}^{\r{iz}} \ni \red{\aves{B \qhatpara{z}}}
\nonumber\\[3mm]&&{}
%
  + \Zi e \frac{\ni}{\mi} I \aves{R^{-2}} \left( - \frac{\red{\psitd}}{q} + \red{\psid} \right)
%
%  + \frac{1}{\mi} \aves{B S_{\r{m}\para}^\r{i}}
%
%  + \frac{1}{\mi} \aves{B S_\r{LC\para}}
\nonumber\\[3mm]&&{}
%
  - \nuL  \ni \red{\aves{B \upara{i}}}
%
  + \nub  \frac{\mb}{\mi} \red{\nb \aves{B \upara{b}}}
%
  - \nun{i} \ni \red{\aves{B \upara{i}}}
%
  - \nuCX \ni \red{\aves{B \upara{i}}}
%
  - \nuLC \ni \red{\aves{B \upara{i}}}
%
  - \nuOL \ni \red{\aves{B \upara{i}}}
\nonumber\\[3mm]&&{}
%
  - \ddV \aves{|\nabla \rho|} v_{\r{i}\zeta} \ni \red{\aves{B \upara{i}}}
%
  + \aves{|\nabla V|} v_{\r{i}\zeta} \ni 
  \biggl(\ddV \frac{\aves{B^2}}{I} \biggl) \frac{I}{\aves{B^2}} \red{\aves{B \upara{i}}}
\nonumber\\[3mm]&&{}
%
  + \ddV \aves{|\nabla V|^2} \chis{i\zeta} \ni
  \frac{\aves{B^2}}{I} \ddV \frac{I}{\aves{B^2}} \red{\aves{B \upara{i}}}
\nonumber\\[3mm]&&{}
%
  - \aves{|\nabla V|^2} \chis{i\zeta} \ni \biggl( \ddV
  \frac{\aves{B^2}}{I} \biggr) \ddV \frac{I}{\aves{B^2}} \red{\aves{B \upara{i}}}
%
  - \frac{1}{\mi} \frac{\aves{B^2}}{I} \ddV \aves{\Pi_\r{i}^\r{res}}
\nonumber\\[3mm]&&{}
%
  - \ddV
    \left[  \aves{|\nabla V|} v_{\r{i}\zeta} \ni
            \frac{\aves{R^2}}{\aves{B^2}}\ave{\frac{|\nabla V|^2}{R^2}}\dd{\psi}{V}(1+2\hat{q}^2)
	    \frac{1}{I} \dd{\psi}{V} BV_{1\r{i}} \right.
\nonumber\\[3mm]&&{}
    \left.- \aves{|\nabla V|^2} \chis{i\zeta} \ni
          \ddV \frac{\aves{R^2}}{\aves{B^2}}\ave{\frac{|\nabla V|^2}{R^2}}\dd{\psi}{V}(1+2\hat{q}^2)
	  \frac{1}{I} \dd{\psi}{V} BV_{1\r{i}}\right]
\\[5mm]
%==============%
% ion toroidal %
%==============%
  \ddt \red{\ni \aves{R \uzt{i}}} &=&
%
  - (\uV{i} - \ugV) \ddV \red{\ni \aves{R \uzt{i}}}
%
  - \dd{\uV{i}}{V} \red{\ni \aves{R \uzt{i}}}
%
  - \ddV \aves{|\nabla V|} v_{\r{i}\zeta} \red{\ni \aves{R \uzt{i}}}
\nonumber\\[3mm]&&{}
  + \ddV \aves{|\nabla V|^2} \chis{i\zeta} \ddV \red{\ni \aves{R \uzt{i}}}
  - \ddV \aves{|\nabla V|^2} \chis{i\zeta} \aves{R \uzt{i}}\ddV \red{\ni} 
  - \frac{1}{\mi} \ddV \aves{\Pi_\r{i}^\r{res}}
\nonumber\\[2.5mm]&&{}
%
  + \hat{\ell}_{11}^\r{ie} \frac{\ni}{\nee} \red{\nee \aves{R \uzt{e}}}
  + \hat{\ell}_{11}^\r{ii}                  \red{\ni  \aves{R \uzt{i}}}
  + \hat{\ell}_{11}^\r{iz} \frac{\ni}{\nz}  \red{\nz  \aves{R \uzt{z}}}
  + \hat{\ell}_{11}^\r{ib} \frac{\ni}{\nb}  \red{\nb  \aves{R \uzt{b}}}
%  + \hat{\ell}_{11}^\r{ib} \frac{I}{\aves{B^2}} \frac{\ni}{\nb} \red{\nb
%  \aves{B \upara{b}}}
\nonumber\\[2.5mm]&&{}
%
%  - \hat{\ell}_{12}^\r{ie} \frac{I}{\aves{B^2}} \ni \red{\aves{B \qhatpara{e}}}
%  - \hat{\ell}_{12}^\r{ii} \frac{I}{\aves{B^2}} \ni \red{\aves{B \qhatpara{i}}}
%  - \hat{\ell}_{12}^\r{iz} \frac{I}{\aves{B^2}} \ni \red{\aves{B \qhatpara{z}}}
  - \hat{\ell}_{12}^\r{ie} \ni \frac{1}{e} \frac{1}{\aves{B^2}} \left( \bri \right)
    \dd{V}{\psi} \ddV \Te
  - \hat{\ell}_{12}^\r{ie} \ni \frac{I}{\aves{B^2}} \red{\aves{B \qhatpara{e}}}
\nonumber\\[2.5mm]&&{}
  + \hat{\ell}_{12}^\r{ii} \ni \frac{1}{\Zi e} \frac{1}{\aves{B^2}} \left( \bri \right)
    \dd{V}{\psi} \ddV \Ti
  - \hat{\ell}_{12}^\r{ii} \ni \frac{I}{\aves{B^2}} \red{\aves{B \qhatpara{i}}}
\nonumber\\[2.5mm]&&{}
  + \hat{\ell}_{12}^\r{iz} \ni \frac{1}{\Zz e} \frac{1}{\aves{B^2}} \left( \bri \right)
    \dd{V}{\psi} \ddV \Tz
  - \hat{\ell}_{12}^\r{iz} \ni \frac{I}{\aves{B^2}} \red{\aves{B \qhatpara{z}}}
\nonumber\\[2.5mm]&&{}
%
  + \Zi \frac{e}{\mi}              \ni \red{\dot{\psi}}
%
  + \Zi \frac{e}{\mi} \dd{\psi}{V} \red{\ni \uV{i}}
%
\nonumber\\[2.5mm]&&{}
%
%  turbulent particle flux model
%
  + \Zi \frac{e^2 \aves{|\nabla V|^2} \De}{\mi\Te} \dd{\psi}{V} 
%%%
%    \left[  \dd{\psi}{V} \frac{\aves{B^2}}{\bri} \frac{\ni+\nb}{\nee}
%            \red{\nee \aves{R \uzt{e}}}
%          - \dd{\psi}{V} \frac{I}{\bri} (\ni+\nb) \red{\aves{B \upara{e}}} \right.
%\nonumber\\[2.5mm]&&{}
%%
%    \left.+ \frac{C_\r{T}-1}{e} \frac{\ni+\nb}{\nee}
%     \left(\ddV \red{\nee \Te} - \Te \ddV \red{\nee} \right)
%	  + (\ni+\nb) \ddV \red{\Phi}
%	  + \frac{\Te}{e} \dd{\psi}{V} C_\r{p} \red{\ni}
%	  + \frac{\Te}{e} \dd{\psi}{V} C_\r{p} \red{\nb} \right]
%%%
  \left[   (\ni + \nb) \dd{\psi}{V} \frac{1}{I} \red{B V_{1\r{e}}}
         + (\ni + \nb) \ddV \red{\Phi} \right.
\nonumber\\[2.5mm]&&{}
  \left. + \frac{C_\r{T}-1}{e} \frac{\ni + \nb}{\nee}
           \left( \ddV \red{\nee \Te} - \Te \ddV \red{\nee} \right)
         + \frac{\Te}{e} \dd{\psi}{V} C_\r{p} \red{\ni}
         + \frac{\Te}{e} \dd{\psi}{V} C_\r{p} \red{\nb} \right]
\nonumber\\[2.5mm]&&{}
%
  - \nuL  \red{\ni \aves{R \uzt{i}}}
%
  + \nub  \frac{\mb}{\mi} \red{\nb \aves{B \uzt{b}}}
%  + \nub  \frac{\mb}{\mi} \frac{I}{\aves{B^2}} \red{\nb \aves{B \upara{b}}}
%
  - \nun{i} \red{\ni \aves{R \uzt{i}}}
%
  - \nuCX \red{\ni \aves{R \uzt{i}}}
%
  - \nuOL \red{\ni \aves{R \uzt{i}}}
\\[5mm]
%==========%
% \ni \Ti  %
%==========%
  \frac{3}{2} \ddt \red{\ni \Ti}  &=& 
%
  - \frac{5}{2} ( \uV{i} - \ugV ) \ddV \red{\ni \Ti}
%
  - \frac{5}{2} \dd{\uV{i}}{V} \red{\ni \Ti}
%
  - \ddV V_\r{i}^\r{hp} \aves{|\nabla V|} \red{\ni \Ti}
\nonumber\\[2.5mm]&&{}
  + \ddV \chis{i} \aves{|\nabla V|^2} \ddV \red{\ni \Ti}
  - \ddV \chis{i} \aves{|\nabla V|^2} \Ti \ddV \red{\ni}
\nonumber\\[2.5mm]&&{}
%
  + ( \uV{i} - \ugV ) \ddV \red{\ni \Ti}
%
  + \uhatth{i} \aves{\vec{B}\cdot\nabla\cdot\tensor{\pi}_\r{i}}
%
  - \frac{3}{2} \nu_\r{Tei} \left(\frac{\nee}{\ni} \red{\ni \Ti} - \red{\nee \Te} \right)
  - \frac{3}{2} \nu_\r{Tiz} \left(\red{\ni \Ti} - \frac{\ni}{\nz}\red{\nz  \Tz} \right)
\nonumber\\[2.5mm]&&{}
%
  - \nuL \left(\red{\ni \Ti}  - \nsup{i} \Ti \right)
%
  - \frac{3}{2} \nuLT{i} (\red{\ni \Ti} - \ni \Tsup{i})
%
\nonumber\\[2.5mm]&&{}
  + \frac{3}{2} \nuion\, \frac{\nee}{\nna + \nnb + \nnc}\, (\red{\nna} \Tna +
  \red{\nnb} \Tnb + \red{\nnc} \Tnc)
\nonumber\\[2.5mm]&&{}
%
  - \frac{3}{2} \nuCX (\red{\ni \Ti} - \red{\ni} \Tna)
%
  + (1 - \fei) \Pb
%
  + \mb \frac{1}{\nb} \frac{\Pb}{\Eb} \frac{\aves{B
  \upara{i}}}{\aves{B^2}} \red{\nb \aves{B \upara{b}}}
%
  + \Palp{i}
%
  + \PRF{i}
\\[5mm]
%=====================%
% 2<B \qipara>/(5 pi) %
%=====================%
 \frac{5}{2} \ni \ddt \red{\aves{B \qhatpara{i}}} &=&
%
  - \frac{1}{\mi} \hat{\mu}_2^\r{i} \red{\aves{B \upara{i}}}
  - \frac{1}{\mi} \hat{\mu}_3^\r{i} \red{\aves{B \qhatpara{i}}}
%
  + \frac{1}{\mi} \hat{\mu}_2^\r{i} \red{B V_{1\r{i}}} 
\nonumber\\[3mm]&&{}
%
%  + \frac{1}{\mi} \hat{\mu}_3^\r{i} {B V_{2\r{i}}}
  - \frac{1}{\Zi e\mi} \hat{\mu}_3^\r{i} I \dd{V}{\psi} \frac{1}{\ni}
  \ddV \red{\ni\Ti}
  + \frac{1}{\Zi e\mi} \hat{\mu}_3^\r{i} I \dd{V}{\psi} \frac{\Ti}{\ni}
  \ddV \red{\ni}
\nonumber\\[3mm]&&{}
%
  - \hat{\ell}_{21}^{\r{ie}} \ni \red{\aves{B \upara{e}}} 
  - \hat{\ell}_{21}^{\r{ii}} \ni \red{\aves{B \upara{i}}} 
  - \hat{\ell}_{21}^{\r{iz}} \ni \red{\aves{B \upara{z}}} 
  - \hat{\ell}_{21}^{\r{ib}} \frac{\ni}{\nb} \red{\nb \aves{B \upara{b}}} 
\nonumber\\[3mm]&&{}
%
  + \hat{\ell}_{22}^{\r{ie}} \ni \red{\aves{B \qhatpara{e}}} 
  + \hat{\ell}_{22}^{\r{ii}} \ni \red{\aves{B \qhatpara{i}}} 
  + \hat{\ell}_{22}^{\r{iz}} \ni \red{\aves{B \qhatpara{z}}} 
\nonumber\\[3mm]&&{}
%
  - \nuL \ni \red{\aves{B \qhatpara{i}}}
%
  - \nun{i} \ni \red{\aves{B \qhatpara{i}}}
\\[5mm]
%=======================%
% <uizt/R> and <R uizt> %
%=======================%
  0 &=&
%
  - \red{\ni \ave{\frac{\uzt{i}}{R}}}
%
  + \frac{1}{\aves{R^2}} \red{\ni \aves{R\uzt{i}}}
%
  - \frac{2\hat{q}^2}{1+2\hat{q}^2} \frac{1}{\aves{R^2}}\red{\ni\aves{R\uzt{i}}}
%
  + \frac{2\hat{q}^2}{1+2\hat{q}^2} \frac{1}{I}\ni \red{\aves{B\upara{i}}}
\\[5mm]
%=======%
% BV_1i %
%=======%
    0 &=&
%
  - \ni \red{BV_{1\r{i}}}
%
  - \frac{1}{\Zi e} \dd{V}{\psi} I \ddV \red{\ni\Ti}
%
  - \ni \dd{V}{\psi} I \ddV \red{\Phi}
\end{eqnarray}
%

%\clearpage

%
\begin{eqnarray}
%======%
% \nz  %
%======%
  \ddt \red{\nz} &=&
%
  - \ddV \red{\nz\uV{z}}
%
  + \ugV \ddV \red{\nz}
%
  + \nee \red{\nnz} f_5S_5 - \nee \red{\nz} A_6 
%
  - \nuL (\red{\nz} - \nsup{i}) 
%
  + S_\r{LC}
%
  - \nuOL \red{\nz}
%
\\[5mm]
%=========%
% radial  %
%=========%
  0 &=&
%
    \Zz \frac{e}{\mz} \nz \left( \bri \right) \frac{1}{I} \red{B V_{1\r{z}}}
%
  + \Zz \frac{e}{\mz} I \nz \red{\aves{B\upara{z}}}
%
  - \Zz \frac{e}{\mz} \aves{B^2} \red{\nz \aves{R\uzt{z}}}
\\[5mm]
%===============%
% <B \upara{z}> %
%===============%
  \nz \ddt \red{\aves{B \upara{z}}} &=&
%
  - \frac{1}{\mi} \hat{\mu}_1^\r{z} \red{\aves{B \upara{z}}}
  - \frac{1}{\mi} \hat{\mu}_2^\r{z} \red{\aves{B \qhatpara{z}}}
%
  + \frac{1}{\mz} \hat{\mu}_1^\r{i} \red{B V_{1\r{z}}} 
\nonumber\\[3mm]&&{}
%
%  + \frac{1}{\mz} \hat{\mu}_2^\r{z} {B V_{2\r{z}}}
  - \frac{1}{\Zz e\mz} \hat{\mu}_2^\r{z} I \dd{V}{\psi} \frac{1}{\nz}
  \ddV \red{\nz\Tz}
  + \frac{1}{\Zz e\mz} \hat{\mu}_2^\r{z} I \dd{V}{\psi} \frac{\Tz}{\nz}
  \ddV \red{\nz}
\nonumber\\[3mm]&&{}
%
  + \hat{\ell}_{11}^{\r{ze}} \nz \red{\aves{B \upara{e}}}
  + \hat{\ell}_{11}^{\r{zi}} \nz \red{\aves{B \upara{i}}}
  + \hat{\ell}_{11}^{\r{zz}} \nz \red{\aves{B \upara{z}}}
  + \hat{\ell}_{11}^{\r{zb}} \frac{\nz}{\nb} \red{\nb \aves{B \upara{b}}}
\nonumber\\[3mm]&&{}
%
  - \hat{\ell}_{12}^{\r{ze}} \nz \red{\aves{B \qhatpara{e}}}
  - \hat{\ell}_{12}^{\r{zi}} \nz \red{\aves{B \qhatpara{i}}}
  - \hat{\ell}_{12}^{\r{zz}} \nz \red{\aves{B \qhatpara{z}}}
\nonumber\\[3mm]&&{}
%
  + \Zz e \frac{\nz}{\mz} I \aves{R^{-2}} \left( - \frac{\red{\psitd}}{q} + \red{\psid} \right)
%
%  + \frac{1}{\mz} \aves{B S_{\r{m}\para}^\r{z}}
%
%  + \frac{1}{\mz} \aves{B S_\r{LC\para}}
\nonumber\\[3mm]&&{}
%
  - \nuL  \nz \red{\aves{B \upara{z}}}
%
  - \nun{z} \nz \red{\aves{B \upara{z}}}
%
  - \nuOL \nz \red{\aves{B \upara{z}}}
\nonumber\\[3mm]&&{}
%
  - \ddV \aves{|\nabla \rho|} v_{\r{z}\zeta} \nz \red{\aves{B \upara{z}}}
%
  + \aves{|\nabla V|} v_{\r{z}\zeta} \nz 
  \biggl(\ddV \frac{\aves{B^2}}{I} \biggl) \frac{I}{\aves{B^2}} \red{\aves{B \upara{z}}}
\nonumber\\[3mm]&&{}
%
  + \ddV \aves{|\nabla V|^2} \chis{z\zeta} \nz
  \frac{\aves{B^2}}{I} \ddV \frac{I}{\aves{B^2}} \red{\aves{B \upara{z}}}
\nonumber\\[3mm]&&{}
%
  - \aves{|\nabla V|^2} \chis{z\zeta} \nz \biggl( \ddV
  \frac{\aves{B^2}}{I} \biggr) \ddV \frac{I}{\aves{B^2}} \red{\aves{B \upara{z}}}
%
  - \frac{1}{\mz} \frac{\aves{B^2}}{I} \ddV \aves{\Pi_\r{z}^\r{res}}
\nonumber\\[3mm]&&{}
%
  - \ddV
    \left[  \aves{|\nabla V|} v_{\r{z}\zeta} \nz
            \frac{\aves{R^2}}{\aves{B^2}}\ave{\frac{|\nabla V|^2}{R^2}}\dd{\psi}{V}(1+2\hat{q}^2)
	    \frac{1}{I} \dd{\psi}{V} BV_{1\r{z}} \right.
\nonumber\\[3mm]&&{}
    \left.- \aves{|\nabla V|^2} \chis{z\zeta} \nz
          \ddV \frac{\aves{R^2}}{\aves{B^2}}\ave{\frac{|\nabla V|^2}{R^2}}\dd{\psi}{V}(1+2\hat{q}^2)
	  \frac{1}{I} \dd{\psi}{V} BV_{1\r{z}}\right]
\\[5mm]
%==============%
% imp toroidal %
%==============%
  \ddt \red{\nz \aves{R \uzt{z}}} &=&
%
  - (\uV{z} - \ugV) \ddV \red{\nz \aves{R \uzt{z}}}
%
  - \dd{\uV{z}}{V} \red{\nz \aves{R \uzt{z}}}
%
  - \ddV \aves{|\nabla V|} v_{\r{z}\zeta} \red{\nz \aves{R \uzt{z}}}
\nonumber\\[3mm]&&{}
  + \ddV \aves{|\nabla V|^2} \chis{z\zeta} \ddV \red{\nz \aves{R \uzt{z}}}
  - \ddV \aves{|\nabla V|^2} \chis{z\zeta} \aves{R \uzt{z}}\ddV \red{\nz} 
  - \frac{1}{\mz} \ddV \aves{\Pi_\r{z}^\r{res}}
\nonumber\\[2.5mm]&&{}
%
  + \hat{\ell}_{11}^\r{ze} \frac{\nz}{\nee} \red{\nee \aves{R \uzt{e}}}
  + \hat{\ell}_{11}^\r{zi} \frac{\nz}{\ni}  \red{\ni  \aves{R \uzt{i}}}
  + \hat{\ell}_{11}^\r{zz}                  \red{\nz  \aves{R \uzt{z}}}
  + \hat{\ell}_{11}^\r{zb} \frac{\nz}{\nb}  \red{\nb  \aves{R \uzt{b}}}
%  + \hat{\ell}_{11}^\r{zb} \frac{I}{\aves{B^2}} \frac{\nz}{\nb} \red{\nb
%  \aves{B \upara{b}}}
\nonumber\\[2.5mm]&&{}
%
%  - \hat{\ell}_{12}^\r{ze} \frac{I}{\aves{B^2}} \nz \red{\aves{B \qhatpara{e}}}
%  - \hat{\ell}_{12}^\r{zi} \frac{I}{\aves{B^2}} \nz \red{\aves{B \qhatpara{i}}}
%  - \hat{\ell}_{12}^\r{zz} \frac{I}{\aves{B^2}} \nz \red{\aves{B \qhatpara{z}}}
  - \hat{\ell}_{12}^\r{ze} \nz \frac{1}{e} \frac{1}{\aves{B^2}} \left( \bri \right)
    \dd{V}{\psi} \ddV \Te
  - \hat{\ell}_{12}^\r{ze} \nz \frac{I}{\aves{B^2}} \red{\aves{B \qhatpara{e}}}
\nonumber\\[2.5mm]&&{}
  + \hat{\ell}_{12}^\r{zi} \nz \frac{1}{\Zi e} \frac{1}{\aves{B^2}} \left( \bri \right)
    \dd{V}{\psi} \ddV \Ti
  - \hat{\ell}_{12}^\r{zi} \nz \frac{I}{\aves{B^2}} \red{\aves{B \qhatpara{i}}}
\nonumber\\[2.5mm]&&{}
  + \hat{\ell}_{12}^\r{zz} \nz \frac{1}{\Zz e} \frac{1}{\aves{B^2}} \left( \bri \right)
    \dd{V}{\psi} \ddV \Tz
  - \hat{\ell}_{12}^\r{zz} \nz \frac{I}{\aves{B^2}} \red{\aves{B \qhatpara{z}}}
\nonumber\\[2.5mm]&&{}
%
  + \Zz \frac{e}{\mz}              \nz \red{\dot{\psi}}
%
  + \Zz \frac{e}{\mz} \dd{\psi}{V} \red{\nz \uV{z}}
%
\nonumber\\[2.5mm]&&{}
%
%  turbulent particle flux model
%
  + \Zz \frac{e \aves{|\nabla V|^2} \De}{\mz\Te} \dd{\psi}{V} 
%%%
%    \left[  (1-C_\r{T})\Te \frac{\nz}{\nee} \ddV \red{\nee}
%           + C_\r{T} \frac{\nz}{\nee} \ddV \red{\nee\Te}
%	  + e \frac{\Te}{e} \dd{\psi}{V} C_\r{p} \red{\nz} \right]
%%%
  \left[   \nz \dd{\psi}{V} \frac{1}{I} \red{B V_{1\r{e}}}
         + \nz \ddV \red{\Phi} \right.
\nonumber\\[2.5mm]&&{}
  \left. + \frac{C_\r{T}-1}{e} \frac{\nz}{\nee}
           \left( \ddV \red{\nee \Te} - \Te \ddV \red{\nee} \right)
         + \frac{\Te}{e} \dd{\psi}{V} C_\r{p} \red{\nz} \right]
\nonumber\\[2.5mm]&&{}
%
  - \nuL  \red{\nz \aves{R \uzt{z}}}
%
  - \nun{z} \red{\nz \aves{R \uzt{z}}}
%
  - \nuOL \red{\nz \aves{R \uzt{z}}}
\\[5mm]
%==========%
% \nz \Tz  %
%==========%
  \frac{3}{2} \ddt \red{\nz \Tz}  &=& 
%
  - \frac{5}{2} ( \uV{z} - \ugV ) \ddV \red{\nz \Tz}
%
  - \frac{5}{2} \dd{\uV{z}}{V} \red{\nz \Tz}
%
  - \ddV V_\r{z}^\r{hp} \aves{|\nabla V|} \red{\nz \Tz}
\nonumber\\[2.5mm]&&{}
  + \ddV \chis{z} \aves{|\nabla V|^2} \ddV \red{\nz \Tz}
  - \ddV \chis{z} \aves{|\nabla V|^2} \Tz \ddV \red{\nz}
\nonumber\\[2.5mm]&&{}
%
  + ( \uV{z} - \ugV ) \ddV \red{\nz \Tz}
%
  + \uhatth{z} \aves{\vec{B}\cdot\nabla\cdot\tensor{\pi}_\r{z}}
%
  - \frac{3}{2} \nu_\r{Tez} \left(\frac{\nee}{\nz} \red{\nz \Tz} - \red{\nee \Te} \right)
  - \frac{3}{2} \nu_\r{Tiz} \left(\frac{\ni }{\nz} \red{\nz \Tz} - \red{\ni  \Ti} \right)
\nonumber\\[2.5mm]&&{}
%
  - \nuL \left(\red{\nz \Tz}  - \nsup{z} \Tz \right)
%
  - \frac{3}{2} \nuLT{z} (\red{\nz \Tz} - \nz \Tsup{z})
%
  + \Palp{z}
%
  + \PRF{z}
\\[5mm]
%=====================%
% 2<B \qzpara>/(5 pz) %
%=====================%
 \frac{5}{2} \nz \ddt \red{\aves{B \qhatpara{z}}} &=&
%
  - \frac{1}{\mz} \hat{\mu}_2^\r{z} \red{\aves{B \upara{z}}}
  - \frac{1}{\mz} \hat{\mu}_3^\r{z} \red{\aves{B \qhatpara{z}}}
%
  + \frac{1}{\mz} \hat{\mu}_2^\r{z} \red{B V_{1\r{z}}}
\nonumber\\[3mm]&&{}
%
%  + \frac{1}{\mz} \hat{\mu}_3^\r{z} {B V_{2\r{z}}}
  - \frac{1}{\Zz e\mz} \hat{\mu}_3^\r{z} I \dd{V}{\psi} \frac{1}{\nz}
  \ddV \red{\nz\Tz}
  + \frac{1}{\Zz e\mz} \hat{\mu}_3^\r{z} I \dd{V}{\psi} \frac{\Tz}{\nz}
  \ddV \red{\nz}
\nonumber\\[3mm]&&{}
%
  - \hat{\ell}_{21}^{\r{ze}} \nz \red{\aves{B \upara{e}}} 
  - \hat{\ell}_{21}^{\r{zi}} \nz \red{\aves{B \upara{i}}} 
  - \hat{\ell}_{21}^{\r{zz}} \nz \red{\aves{B \upara{z}}} 
  - \hat{\ell}_{21}^{\r{zb}} \frac{\nz}{\nb} \red{\nb \aves{B \upara{b}}} 
\nonumber\\[3mm]&&{}
%
  + \hat{\ell}_{22}^{\r{ze}} \nz \red{\aves{B \qhatpara{e}}} 
  + \hat{\ell}_{22}^{\r{zi}} \nz \red{\aves{B \qhatpara{i}}} 
  + \hat{\ell}_{22}^{\r{zz}} \nz \red{\aves{B \qhatpara{z}}} 
\nonumber\\[3mm]&&{}
%
  - \nuL \nz \red{\aves{B \qhatpara{z}}}
%
  - \nun{z} \nz \red{\aves{B \qhatpara{z}}}
\\[5mm]
%=======================%
% <uzzt/R> and <R uzzt> %
%=======================%
  0 &=&
%
  - \red{\nz \ave{\frac{\uzt{z}}{R}}}
%
  + \frac{1}{\aves{R^2}} \red{\nz \aves{R\uzt{z}}}
%
  - \frac{2\hat{q}^2}{1+2\hat{q}^2} \frac{1}{\aves{R^2}}\red{\nz\aves{R\uzt{z}}}
%
  + \frac{2\hat{q}^2}{1+2\hat{q}^2} \frac{1}{I}\nz \red{\aves{B\upara{z}}}
\\[5mm]
%=======%
% BV_1z %
%=======%
    0 &=&
%
  - \nz \red{BV_{1\r{z}}}
%
  - \frac{1}{\Zz e} \dd{V}{\psi} I \ddV \red{\nz\Tz}
%
  - \nz \dd{V}{\psi} I \ddV \red{\Phi}
\end{eqnarray}
%

\clearpage

%
\begin{eqnarray}
%======%
% \nb  %
%======%
  \ddt \red{\nb} &=&
%
  - \ddV \red{\nb\uV{b}}
%
  + \ugV \ddV \red{\nb}
%
  +  \frac{\Pb}{\Eb}
%
  - g \frac{\Pbperp}{\Eb}
%
  - \nub \red{\nb}
%
  - \nubL \red{\nb}
%
  - \nueff g (\sqrt{\delta} \red{\nb} - \red{\nbrp})
%
\nonumber\\[2.5mm]&&{}
  + \ddV \aves{|\nabla V|^2} f_t \Db \ddV \red{\nb}
%
\\[5mm]
%=========%
% radial  %
%=========%
  0 &=&
%
    \Zb \frac{e}{\mb} \left( \bri \right) \frac{1}{I} \red{\nb B V_{1\r{b}}}
%
  + \Zb \frac{e}{\mb} I \red{\nb \aves{B\upara{b}}}
%
  - \Zb \frac{e}{\mb} \aves{B^2} \red{\nb \aves{R\uzt{b}}}
\\[5mm]
%===============%
% \nb \upara{b} %
%===============%
  \ddt \red{\nb \aves{B \upara{b}}} &=&
%
  - \frac{1}{\mb \nb} \hat{\mu}_1^\r{b} \red{\nb \aves{B \upara{b}}}
%
  + \frac{1}{\mb \nb} \hat{\mu}_1^\r{b} \red{\nb B V_{1\r{b}}}
\nonumber\\[2.5mm]&&{}
%
  + \hat{\ell}_{11}^\r{bb} \frac{\me\nee}{\mb\nb} \red{\nb \aves{B \upara{b}}}
%
  + \sum_{a\neq\r{b}} \hat{\ell}_{11}^{\r{b}a} \frac{\me\nee}{\mb} \red{\aves{B u_{a\para}}}
\nonumber\\[2.5mm]&&{}
%
  + \Zb \frac{e}{\mb} \nb I \aves{R^{-2}} \left( - \frac{\red{\psitd}}{q} + \red{\psid} \right)
%
  - \nubL \red{\nb \aves{B\upara{b}}}
%
  - \nub  \red{\nb \aves{B\upara{b}}}
%
  - \nun{b} \red{\nb \aves{B\upara{b}}}
\nonumber\\[2.5mm]&&{}
%
  - \nuCX \red{\nb \aves{B\upara{b}}}
%
%  + \vbph \frac{\Pbpara}{\Eb}
%
  + \frac{1}{\mb} \aves{B S_{\r{m}\para\r{b}}}
%
\\[5mm]
%===============%
% beam toroidal %
%===============%
  \ddt \red{\nb \aves{R \uzt{b}}} &=&
%
  \ugV \ddV \red{\nb \aves{R \uzt{b}}}
%
  - \ddV \aves{R \uzt{b}} \red{\nb \uV{b}} 
%
%  - \ddV [\aves{|\nabla V|} v_{\r{b}\zeta} \red{\nb \aves{R \uzt{b}}}]
%\nonumber\\[3mm]&&{}
%  + \ddV \aves{|\nabla V|^2} \chis{b\zeta} \ddV \red{\nb \aves{R \uzt{b}}}
%  - \ddV \aves{|\nabla V|^2} \chis{b\zeta} \aves{R \uzt{b}}\ddV \red{\nb} 
%  - \frac{1}{\mb} \ddV \aves{\Pi_\r{b}^\r{res}}
\nonumber\\[2.5mm]&&{}
%
  + \frac{1}{\mb} \ell_{11}^\r{be} \frac{1}{\nee} \red{\nee \aves{R \uzt{e}}}
  + \frac{1}{\mb} \ell_{11}^\r{bi} \frac{1}{\ni}  \red{\ni  \aves{R \uzt{i}}}
  + \frac{1}{\mb} \ell_{11}^\r{bz} \frac{1}{\nz}  \red{\nz  \aves{R \uzt{z}}}
%  + \hat{\ell}_{11}^\r{bb} \red{\ni \aves{R \uzt{b}}}
  + \frac{1}{\mb} \ell_{11}^\r{bb} \frac{1}{\nb}  \red{\nb  \aves{R \uzt{b}}}
\nonumber\\[2.5mm]&&{}
%
%  - \frac{1}{\mb} \ell_{12}^\r{be} \frac{I}{\aves{B^2}} \red{\aves{B \qhatpara{e}}}
%  - \frac{1}{\mb} \ell_{12}^\r{bi} \frac{I}{\aves{B^2}} \red{\aves{B \qhatpara{i}}}
%
  + \Zb \frac{e}{\mb}              \nb \red{\dot{\psi}}
%
  + \Zb \frac{e}{\mb} \dd{\psi}{V} \red{\nb \uV{b}}
%
\nonumber\\[2.5mm]&&{}
%%
%%  turbulent particle flux model
%%
%  + \Zi \frac{e^2 \aves{|\nabla V|^2} \De}{\mi\Te} \dd{\psi}{V} 
%%
%    \left[  \dd{\psi}{V} \frac{\aves{B^2}}{\bri} \red{\nee \aves{R \uzt{e}}}
%          - \dd{\psi}{V} \frac{I}{\bri} \nee \red{\aves{B \upara{e}}} \right.
%\nonumber\\[2.5mm]&&{}
%%
%    \left.+ \frac{C_\r{T}-1}{e} \left(\ddV \red{\nee \Te} - \Te \ddV \red{\nee} \right)
%	  + \nee \ddV \red{\Phi}
%	  + \frac{\Te}{e} \dd{\psi}{V} C_\r{p} \red{\nee} \right]
%\nonumber\\[2.5mm]&&{}
%
  - \nubL  \red{\nb \aves{R \uzt{b}}}
%
  - \nub  \red{\nb \aves{B \uzt{b}}}
%
  - \nun{i} \red{\nb \aves{R \uzt{b}}}
%
  - \nuCX \red{\nb \aves{R \uzt{b}}}
%
%  - \nuOL \red{\nb \aves{R \uzt{b}}}
%
  + \frac{1}{\mb} \frac{I}{\aves{B^2}} \aves{B S_{\r{m}\para\r{b}}}
%
\\[5mm]
%=======================%
% <ubzt/R> and <R ubzt> %
%=======================%
  0 &=&
%
  - \red{\nb \ave{\frac{\uzt{b}}{R}}}
%
  + \frac{1}{\aves{R^2}} \red{\nb \aves{R\uzt{b}}}
%
  - \frac{2\hat{q}^2}{1+2\hat{q}^2} \frac{1}{\aves{R^2}}\red{\nb\aves{R\uzt{b}}}
%
  + \frac{2\hat{q}^2}{1+2\hat{q}^2} \frac{1}{I}\red{\nb \aves{B\upara{b}}}
\\[5mm]
%=======%
% BV_1b %
%=======%
    0 &=&
%
  - \red{\nb BV_{1\r{b}}}
%
  - \frac{1}{\Zb e} \dd{V}{\psi} I \red{\nb} \ddV \Tb 
%
  - \frac{1}{\Zb e} \dd{V}{\psi} I \Tb \ddV \red{\nb}
%
  - \nb \dd{V}{\psi} I \ddV \red{\Phi}
\\[5mm]
%=======%
% \nna  %
%=======%
  \ddt \red{\nna} &=& 
%
  - \dd{\ugV}{V} \red{\nna}
%
  + \ddV \aves{|\nabla V|^2} D_{01} \ddV \red{\nna}
%
  - \nuion\, \frac{\nee}{\nna + \nnb + \nnc}\, \red{\nna}
%
\nonumber\\[2.5mm]&&{}
  - \nuCX\, \frac{\ni}{\nna + \nnb}\, \red{\nna}
%
  + \gamma_0 \nuL (\red{\ni} - \nsup{i})
%
  + \ddV V' \Gamma^\r{puff}\biggr|_\r{b}
\\[5mm]
%=======%
% \nnb  %
%=======%
  \ddt \red{\nnb} &=& 
%
  - \dd{\ugV}{V} \red{\nnb}
%
  + \ddV \aves{|\nabla V|^2} D_{02} \ddV \red{\nnb}
%
  - \nuion\, \frac{\nee}{\nna + \nnb + \nnc}\, \red{\nnb}
%
%\nonumber\\[2.5mm]&&{}
  + \nuCX\, \frac{\ni}{\nna + \nnb}\, \red{\nna}
%
%  + \fCX \frac{\Pb}{\Eb}
\\[5mm]
%=======%
% \nnc  %
%=======%
  \ddt \red{\nnc} &=& 
%
  - \dd{\ugV}{V} \red{\nnc}
%
  + \ddV \aves{|\nabla V|^2} D_{03} \ddV \red{\nnc}
%
  - \nuion\, \frac{\nee}{\nna + \nnb + \nnc}\, \red{\nnc}
%
  + \fCX \frac{\Pb}{\Eb}
\\[5mm]
%=======%
% \nnz  %
%=======%
  \ddt \red{\nnz} &=& 
%
  - \dd{\ugV}{V} \red{\nnz}
%
  + \ddV \aves{|\nabla V|^2} D_{n\r{z}} \ddV \red{\nnz}
%
  + \ddV \aves{|\nabla V|^2} \Delta D_{\r{C0}} \ddV n_\r{C0}
%
  - \nee \red{\nnz} f_5S_5 + \nee \red{\nz} A_6 
\nonumber\\[2.5mm]&&{}
%
  + \gamma_\r{0z} \nuL \, (\red{\nz} - \nsup{z})
%
  + \gamma_\r{i,z} \nuL \, (\red{\ni} - \nsup{i})
%
  - (1-\gamma_\r{0z}) \nuL \, (\red{\nnz} - n_\r{C0})
\nonumber\\[2.5mm]&&{}
%
  + \ddV V' \Gamma^\r{puff}_\r{z}\biggr|_\r{b}
\\[5mm]
%========%
% \nbrp  %
%========%
  \ddt (g \red{\nbrp}) &=&
%
    g \frac{\Pbperp}{\Eb}
%
  + g \nueff (\sqrt{\delta} \red{\nb} - \red{\nbrp})
%
  - g \nub \red{\nbrp}
%
\nonumber\\[2.5mm]&&{}
  - \ddV g \aves{|\nabla V|} \ubrp \red{\nbrp}
  + \ddV g \aves{|\nabla V|^2} \Dbrp \ddV \red{\nbrp}
\end{eqnarray}

\bigskip

\section{Boundary Conditions}

%
\begin{tabular}{|c|c|c|c|c|c|}\hline
 num. & name & variable               & \shortstack{highest order\\ derivative} & center & edge \\\hline
 1    & LQm1 & $\Phi$                 & 2    &  N           & 0\\
 2    & LQm2 & $\psitd$               & 0    &  $\times$    & $\times$\\
 3    & LQm3 & $\psid$                & 0    &  $\times$    & $\times$\\
 4    & LQm4 & $\psi$                 & 2    &  N           & $2\pi\mu_0I_\r{p}$\\
 5    & LQm5 & $\psit/\mu_0$          & 2    &  0           & $\aves{|\nabla V|^2/R^2}|_{V_b}I_\r{vac}/(4\pi^2)$\\
 6    & LQe1 & $\nee$                 & 2    &  N           & N$^\dagger$\\
 7    & LQe2 & $\nee\uV{e}$           & 1    &  0           & $\times$\\
 8    & LQe3 & $\aves{B\upara{e}}$    & 2(0) &  N($\times$) & N($\times$)\\
 9    & LQe4 & $\nee\aves{R\uzt{e}}$  & 2    &  N           & N\\
 10   & LQe5 & $\nee \Te$             & 2    &  N           & N\\
 11   & LQe6 & $\aves{B\qhatpara{e}}$ & 0    &  $\times$    & $\times$\\
 12   & LQe7 & $\nee\aves{\uzt{e}/R}$ & 0    &  $\times$    & $\times$\\
 13   & LQe8 & $B V_{1\r{e}}$         & 0    &  $\times$    & $\times$\\
 14   & LQi1 & $\ni$                  & 2    &  N           & N$^\dagger$\\
 15   & LQi2 & $\ni\uV{i}$            & 1    &  0           & $\times$\\
 16   & LQi3 & $\aves{B\upara{i}}$    & 2(0) &  N($\times$) & N($\times$)\\
 17   & LQi4 & $\ni\aves{R\uzt{i}}$   & 2    &  N           & N\\
 18   & LQi5 & $\ni \Ti$              & 2    &  N           & N\\
 19   & LQi6 & $\aves{B\qhatpara{i}}$ & 0    &  $\times$    & $\times$\\
 20   & LQi7 & $\ni\aves{\uzt{i}/R}$  & 0    &  $\times$    & $\times$\\
 21   & LQi8 & $B V_{1\r{i}}$         & 0    &  $\times$    & $\times$\\
 22   & LQz1 & $\nz$                  & 2    &  N           & N$^\dagger$\\
 23   & LQz2 & $\nz\uV{z}$            & 1    &  0           & $\times$\\
 24   & LQz3 & $\aves{B\upara{z}}$    & 2(0) &  N($\times$) & N($\times$)\\
 25   & LQz4 & $\nz\aves{R\uzt{z}}$   & 2    &  N           & N\\
 26   & LQz5 & $\nz \Tz$              & 2    &  N           & N\\
 27   & LQz6 & $\aves{B\qhatpara{z}}$ & 0    &  $\times$    & $\times$\\
 28   & LQz7 & $\nz\aves{\uzt{z}/R}$  & 0    &  $\times$    & $\times$\\
 29   & LQz8 & $B V_{1\r{z}}$         & 0    &  $\times$    & $\times$\\
 30   & LQb1 & $\nb$                  & 0(2) &  $\times$(N) & $\times$(N)\\
 31   & LQb2 & $\nb\uV{b}$            & 1    &  0           & 0$^\ddagger$\\
 32   & LQb3 & $\nb\aves{B\upara{b}}$ & 0    &  $\times$    & $\times$\\
 33   & LQb4 & $\nb\aves{R\uzt{b}}$   & 0    &  $\times$    & $\times$\\
 34   & LQb7 & $\nb\aves{\uzt{b}/R}$  & 0    &  $\times$    & $\times$\\
 35   & LQb8 & $B V_{1\r{b}}$         & 0    &  $\times$    & $\times$\\
 36   & LQn1 & $\nna$                 & 2    &  N           & $\aves{|\nabla V|}S_\r{gas}|_b$\\
 37   & LQn2 & $\nnb$                 & 2    &  N           & 0\\
 38   & LQn3 & $\nnc$                 & 2    &  N           & 0\\
 39   & LQnz & $\nnz$                 & 2    &  N           & $\aves{|\nabla V|}S_\r{gas,z}|_b$\\
 40   & LQr1 & $\nbrp$                & 0(2) &  $\times$(N) & $\times$(N)\\\hline
 \multicolumn{5}{l}{0: Dirichlet condition $(u=0)$,\ 
                    N: Neumann condition $(u'=0)$}\\
 \multicolumn{5}{l}{value: Neumann condition $(u'=\mbox{value})$,\ 
                    $\times$: no condition}\\
 \multicolumn{5}{l}{$\dagger$: imposed as a constraint in V5.60 (\texttt{MDLWBC} = 1), see Sec.~\ref{sec:wallbc}}\\
 \multicolumn{5}{l}{$\ddagger$: V5.61, \texttt{MDBEAM} = 1 only, see Sec.~\ref{sec:beamwall}}
\end{tabular}
%

\medskip

NOTE: The dependent variable of LQm4 turns to be $\psi/\mu_0$ when NB is
applied. Subsequently, the edge boundary condition listed for that variable
is also changed to $2\pi I_\r{p}$.

NOTE: Rows 2 to 5 deserve a word, because the variable that carries the
second-order operator is not the variable of the equation in which the
operator is written. The LQm2 equation advances $\psitd$ and contains
$\ddV[\ave{|\nabla V|^2/R^2}\aves{R^{-2}}^{-1}\ddV(\psit/\mu_0)]$; the
LQm3 equation advances $\psid$ and contains
$\ddV[\ave{|\nabla V|^2/R^2}\ddV\psi]$. Counting per dependent variable
over the whole system, as this section does, $\psitd$ and $\psid$ are never
differentiated and so take order 0 and no condition of their own, while
$\psit/\mu_0$ and $\psi$ take order 2 and carry the two conditions each.
The table above now reads that way.
In the implementation the two conditions of each pair are applied to the
LQm2 and LQm3 equation rows, since that is where the boundary term left by
the integration by parts appears; this is the correct weak form and is not
a discrepancy with the table. The axis condition of $\psit/\mu_0$ is imposed
there as $\psitd|_{V=0}=0$, which holds $\psit(0)$ at its initial value,
namely zero, and is therefore the Dirichlet condition listed.

NOTE: In the table the column labeled by the highest order derivative means
the order of the highest order derivative of a dependent variable, not
an equation, in the corresponding row.
Please note that the term not integrated by parts is never counted.
For example, the first row shows the highest order of derivatives is
equal to 2 for $\Phi$ because the LQm1 equation
has the diffusion term having the second derivative of $\Phi$ with
respect to $V$. In addition, the LQe2 and LQi2 equations also have the
first derivative of $\Phi$, which is not a dependent variable of these
equations. As a whole, the highest order derivative of $\Phi$ in this
equation system is 2 that comes from the LQm1 equation.
When we look at $\nee$, we find that the highest order derivative of
$\nee$ is 2 shown in the LQe5 equation, even though no derivative
appears in the LQe1 equation. The Neumann conditions for $\nee$ come
from the LQe5 equation.
When it comes to $\nee\uV{e}$, the highest derivative is found as order 1
in the LQe1 equation, but this term is not integrated by parts.
It means that any boundary condition is not naturally imposed on
$\nee\uV{e}$. The Dirichlet boundary condition should be imposed at axis.





\clearpage
\section{Momentum pinch of V5.62}
\setcounter{equation}{0}
\label{sec:v562}

V5.62 changes the turbulent momentum pinch of the thermal species.  The flux of
toroidal angular momentum of species $s$ is written with the density outside the
gradient (the form of the scaling used below),
\begin{equation}
  \Gamma_{\r{L},s} = m_s\left(\aves{|\nabla V|}\,V_\r{p}\,n_s\aves{Ru_{s\zeta}}
   - \aves{|\nabla V|^2}\,\mu_s\,n_s\,\partial_V\aves{Ru_{s\zeta}}\right),
\end{equation}
with the momentum diffusivity $\mu_s$ and the pinch velocity $V_\r{p}$ ($V_\r{p}>0$
outward).  The pinch is
\begin{equation}
  V_\r{p} = -\frac{\mu_\r{i}}{R_0}\left[1.1\,g\!\left(\frac{R_0}{L_n}\right) + 1.0\right] T(\rho),
  \qquad \frac{R_0}{L_n} = -R_0\,\frac{d\ln n_\r{e}}{d\rho_V},
  \label{eq:v562pinch}
\end{equation}
where $\rho_V$ is the volume radius and $n_\r{e}$ is smoothed as before.

\paragraph{Origin and range of the scaling.}  The multi-machine scaling of Tala
et al.\ \cite{Tala2012} is $-R_0V_\r{pinch}/\chi_\phi = 1.1\,R/L_n + 1.0$, both terms
inward.  It was obtained from NBI modulation in JET, DIII-D and AUG and from
magnetic perturbations in NSTX, with the pinch and the Prandtl number averaged over
$0.4<\rho_\r{tor}<0.8$; the data cover $R/L_n$ = 1.2--3.6 (0.2--3.8 with NSTX) and
pinch numbers 0.9--6.2.  The slopes of the single devices are 1.2 (JET, DIII-D), 1.5
(AUG) and 0.6 (NSTX); gyrokinetic calculations give a dependence about half as strong.
The profiles are ``not accurate in the edge region'', and an edge torque added to the
analysis does not change the coefficients inside $r/a<0.8$.  The analysis method
\cite{Mantica2010} and the JET subset \cite{Tala2011} write the flux in the form
above; residual stress and intrinsic torque are neglected.  No upper limit of $R/L_n$
is stated.  Beyond the range, the JETPEAK database \cite{Weisen2012} ($R/L_n$ = 1--10,
$r/a$ up to 0.85) gives an $R/L_n$ coefficient of 0.41--0.66 for the steady-state
rotation gradient (pinch and residual stress together), and the regression of GKW
calculations 0.44; neither shows a saturation.

\paragraph{The closure.}  Equation (\ref{eq:v562pinch}) is a reduced closure, not the
Tala law extrapolated over the whole plasma: it is the Tala fit inside its data, it
takes the slope beyond the data from the multi-machine and gyrokinetic information,
and it is not applied in the edge and the SOL.
\begin{itemize}
\item $g(x) = x$ for $|x|\le x_1$ and
\begin{equation}
  g(x) = \mathrm{sgn}(x)\left[x_1 + s(|x|-x_1) + (1-s)\,w\tanh\frac{|x|-x_1}{w}\right]
  \quad (|x|>x_1),
\end{equation}
with $x_1$ = 3.8 (\texttt{PCHRLN}, the largest $R/L_n$ of the fit), $s$ = 0.4
(\texttt{PCHSLP}, so that the slope goes from 1.1 to 0.44) and $w$ = 1 (\texttt{PCHWID}).
$g$ is continuous up to its second derivative.  The pinch number is 7.7 at $R/L_n$ = 8
and 8.6 at 10, within the measured pinch numbers (up to about 8), against 9.8 and 12 for
the linear law.  $s$ = 0 saturates the pinch number at $1+1.1(x_1+w)$; this is an
experimental option without a basis in the literature.
\item $T(\rho) = 1$ for $\rho\le\rho_1$ = 0.8 (\texttt{RhoPch1}, the outer edge of the
data), $0$ for $\rho\ge\rho_2$ = 0.9 (\texttt{RhoPch2}) and the quintic smoothstep
$1-\xi^3(10-15\xi+6\xi^2)$, $\xi=(\rho-\rho_1)/(\rho_2-\rho_1)$, between.  $T$ is
continuous up to its second derivative and exactly zero outside $\rho_2$, so that the
pinch carries no momentum through the wall.  The location follows the data
(0.8), the momentum boundary of the TOPICS validation (0.9 \cite{Honda2013}) and the
pedestal-top boundaries of present integrated modelling (0.8--0.94); the width is a
numerical choice.  Outside $\rho_2$ the momentum is transported by the diffusion only.
Cubic, cosine and tanh tapers with the same end points give the same rotation to three
digits.
\item $V_\r{p}$ is built from the ion viscosity and given to every thermal species; the
beam has no turbulent momentum flux.  With a common $\mu$ the fluxes sum to
$V_\r{p}L - \mu\sum_s m_sn_s\partial_V\aves{Ru_{s\zeta}}$ times the metric factors, the
pinch of the total thermal momentum $L = \sum_s m_sn_s\aves{Ru_{s\zeta}}$ used in
TOPICS \cite{Honda2013}, whether or not the species rotate alike.  With
species-dependent \texttt{rMus0} the pinch number $-R_0V_\r{p}/\mu_s$ differs between
species and this equivalence is lost.
\end{itemize}

\paragraph{Correction of the sign.}  Up to V5.61 the code evaluated
$V_\r{p} = (\mu/R_0)(1.1R_0\,d\ln n_\r{e}/d\rho_V + 1)$, i.e.\
$-R_0V_\r{p}/\mu = 1.1R_0/L_n - 1$: the constant term pointed outward, against the
scaling it cited.  V5.62 uses $+1.0$ (verified on the output pinch velocity to three
digits).  The commented-out edge suppressions of the old code (a tanh at $\rho$ = 0.98
and a Gaussian at $\rho$ = 1) had no documented basis and are removed.  Without the
radial restriction the old pinch reached a pinch number of 14 at $\rho$ = 0.9 and about
140 at the separatrix, and the convective flux at the wall carried momentum into the
plasma (+0.04 N\,m in the JT-60U case).

\paragraph{Validation.}  JT-60U-like case of Sec.~\ref{sec:v561val} (3.2 MW co-tangential and 0.4 MW
perpendicular NBI, $\chi=\mu=0.22(1+19\rho^2)$ m$^2$/s, \texttt{PNsSUP(1)} = 0.0125 as explained
below), production defaults (\texttt{FSTPTM} = 3), 1.5 s:
\begin{center}
\begin{tabular}{|l|c|c|c|c|c|c|}\hline
case & $u_{\r{i}\phi}(0)$ & $u_{\r{i}\phi}(0.82)$ & $E_r(0.5)$ & $j_r(0.97)$ & $j\times B$ torque & $L_\phi$ \\
 & km/s & km/s & kV/m & A/m$^2$ & N\,m & N\,m\,s \\\hline
no ripple & 510 & 57 & $+44$ & $-0.017$ & $+0.16$ & 2.03 \\
half ripple (FST) & 370 & 27 & $+22$ & $-0.103$ & $-3.25$ & 1.58 \\
full ripple & 198 & $-8$ & $-5$ & $-0.180$ & $-7.28$ & 1.05 \\\hline
\end{tabular}
\end{center}
The line-averaged density is $1.15\times10^{19}$ m$^{-3}$ and the thermal mass inside $\rho$ = 1 is 1.54--1.62 mg.
The first version of this table used \texttt{PNsSUP(1)} = 0.01 (the default) and \texttt{FSTPTM} = 2:
621, 441 and 226 km/s ($\bar n_\r{e} = 1.02\times10^{19}$ m$^{-3}$, thermal mass 1.29 mg);
\texttt{FSTPTM} = 3 gives 620 and 223 km/s without and with full ripple at 0.01, and V5.61 gave 658,
417 and 128 km/s.  The ripple changes the rotation, $E_r$, the return current and its torque at
\texttt{PNsSUP(1)} = 0.0125 as at 0.01.  The rest of this paragraph refers to
\texttt{PNsSUP(1)} = 0.01.  The pinch number is 6.3 at $\rho$ = 0.5 and 7.9 at
0.8 and vanishes outside 0.9; in V5.61 it was 14 at 0.9.  The profile is more peaked
in the core (the constant term) and lower at the edge ($u_{\r{i}\phi}(0.9)$ = 46
against 105 km/s).  The volume integrals close to $2\times10^{-5}$ of the injected
torque in every row, the pinch flux at the wall is zero and the friction terms sum to
$3\times10^{-9}$ N\,m.  With \texttt{NRMAX} = 120 the total angular momentum agrees
with \texttt{NRMAX} = 60 to 0.5\% up to 0.8 s (the Picard iteration is limited there by
the known impurity terms at the axis, as in V5.61).  On the test equilibria of
Sec.~\ref{sec:v561val} densities and temperatures change by less than $3\times10^{-3}$
and the rotation by 13--40\%; the mean Picard count rises by 1--4 and no step more
reaches \texttt{ICMAX}.  The Tala law with the corrected sign but without the radial
restriction and the slope reduction spins the JT-60U case up to 881 km/s and stops at
1.4 s at the impurity grid mode of supersonic rotation.

\paragraph{Density of the benchmark.}  \texttt{PNsSUP(1)} is the upstream SOL density and sets
the density level of the whole plasma (Sec.~\ref{sec:v552other}); the shape of the density profile
does not depend on it.  The benchmark input sets it to 0.0125 ($10^{20}$ m$^{-3}$), which gives
$\bar n_\r{e} = 1.15\times10^{19}$ m$^{-3}$ and $n_\r{e}(\rho=1) = 0.22\times10^{19}$ m$^{-3}$.  The
line-averaged density is the one estimated from the measured profile of the JT-60U L-mode
discharge compared in \cite{NF2008} (Fig.~1 of \cite{Yoshida2006}, which gives no value), and
the separatrix density lies at the lower end of the range expected for JT-60U L-modes
(0.2--0.45$\times10^{19}$ m$^{-3}$).  This value normalizes the density of this benchmark; it is
not a universal value, and the default of \texttt{PNsSUP(1)} stays 0.01.  A new device,
discharge, equilibrium, heating scheme or particle-transport model needs its own value, found
by an outer loop of two or three runs to the target density (secant in $\log$--$\log$; there is
no density feedback inside the code).  The benchmark is JT-60U-like and not the discharge:
$I_\r{p}$ = 1.0 MA against 1.23 MA, a circular cross-section against a D-shaped plasma of
77 m$^3$, and the computed density profile falls more steeply towards the edge
($n_\r{e}(0.9)/\bar n_\r{e}$ = 0.45 against 0.74).  The rotation was not used to choose
\texttt{PNsSUP(1)}: it is a consistency check of the momentum model, not a fitting target.

The reduced ripple model is judged by the changes with the ripple rather than by the
absolute rotation.  The co-rotation without ripple is 1.7 times that of \cite{NF2008}
(510 against 300 km/s); the
states of that work, restored from its save files with the 2008 code, show why.  The
2008 plasma was pure deuterium with a flat density ($n_\r{e}$ = 1.4--1.5$\times10^{19}$
m$^{-3}$ up to $r/a$ = 0.7), 1.4 times heavier than the present one (thermal mass 2.30
against 1.59 mg, moment of inertia $2.70$ against $1.87\times10^{-5}$ kg\,m$^2$); its
torque was 7.63 N\,m (85 keV, the whole beam velocity toroidal at $R_0$) with half of it
deposited off axis, it had no momentum pinch, and its wall held $u_\phi = 0$.  The
momentum confinement times $L_\phi/T_\r{NBI}$ of the thermal species are similar (0.137
and 0.143 s), so the mass-weighted rotation $R_0L_\phi/I_\phi$ differs by 1.3 (132 and 174
km/s) and the more peaked profile gives the rest of the factor 1.7 on axis.

\paragraph{Limitations.}  The pinch number is taken from the electron density gradient
alone; the dependences on $q$, $\hat{s}$, $\varepsilon$, $T_\r{i}/T_\r{e}$ found in
the regressions are not included, and no residual stress is modelled.  The radial
restriction is fixed in $\rho$ and does not follow the pedestal of an H-mode.  The
momentum diffusivity itself is prescribed ($\mu = \chi$ in the benchmark).  A
Coriolis-type alternative with a geometric dependence ($-R_0V/\mu\propto\sqrt{r/R}$, as
used for DTT \cite{Casiraghi2021}) rotates the JT-60U case 30\% more slowly; its
coefficient is a calibration, not a universal constant, and it is not implemented in
the production code.


\section{Beam and ripple model of V5.61}
\setcounter{equation}{0}
\label{sec:v561}

V5.61 changes the beam and ripple equations.  The model itself (its
assumptions, the ordering behind each term and the comparison with the 2008
ripple model \cite{NF2008}) is described in \texttt{beam\_ripple.tex}; this
section states what the code solves and what changed.

\begin{center}
\begin{tabular}{|l|l|}\hline
change & effect \\\hline
ripple momentum terms of the beam removed (B-full) & \texttt{FSRP} = 1 with NBI \\
LQr1 written as the conservation law of $g\nbrp$ & \texttt{FSRP} = 1 with NBI \\
$\Gamma_\r{b}(\rho_b) = 0$ & \texttt{MDBEAM} = 1 with NBI \\
fast-ion friction $\ell_{11}^{\r{f}k} = \ell_{11}^{k\r{f}}$ restored & NBI (\texttt{MDLNEO} = 1, 2) \\
$Z_\r{eff}$ and the fast-ion inventory weighted by $g$ & \texttt{FSRP} = 1 with NBI \\
\texttt{FSRP} restricted to 0, 1, 2; \texttt{FSRP} = 2 keeps $g P_{\r{b}\perp}$ in $\nb$ & input \\
on-axis ripple trapping rate from the interior & \texttt{FSRP} = 1 with NBI \\\hline
\end{tabular}
\end{center}

Runs without NBI are bit-identical to V5.60 (X-point and circular test
equilibria, \texttt{NRMAX} = 60 and 240, \texttt{MDLNZD} = 0, \texttt{MDLN0FL} = 2,
\texttt{FSRP} = 0 and 1).  No default switch changes.  For runs with NBI, and in
particular with \texttt{FSRP} = 1, \texttt{NANDER} = $-20$ is recommended
(Sec.~\ref{sec:anderson}): the impurity radial flux at the wall converges slowly
without it once the beam heats the edge.

\subsection{\texttt{FSRP} and the ripple-well fraction}

\texttt{FSRP} is a mode, not a factor: 0, no ripple; 1, full ripple model
(ripple-trapped ions, LQr1, and ripple diffusion of $\nb$); 2, ripple diffusion
of $\nb$ without ripple-trapped ions.  Other values are rejected.  The ripple
amplitude is $\delta(r,\theta) = \delta_\r{min} I_0(N a_\delta/R_{\delta\r{min}})$
with \texttt{DltRPn} $=\delta_\r{min}$; \texttt{FSRP} only switches it on.
$R_{\delta\r{min}}$ (input \texttt{RDltRPn}, default 2.4 m) is the major radius on the
midplane where the ripple amplitude has its minimum $\delta_\r{min}$; the formula takes
the surfaces $\delta$ = const as circles centred there, a property of the TF coil
geometry (the coil shape and radial position), not of the plasma.  2.4 m is the JT-60U
value that entered the 2008 code together with the ripple model (``specific value for
JT-60U'', also used by the routine that drew the ripple contours); it was fixed in the
code up to V5.61.00, and other devices need their own value.  The text of \cite{NF2008}
gives 2.8 m, but
its own numbers require 2.4 m: with $\delta_\r{min} = 2\times10^{-6}$ and $N = 18$ the
formula gives $\delta_a = 1.28$\% at the outboard midplane and $0.30$\% at
$\theta = \pi/2$ with 2.4 m, the values quoted in \cite{NF2008}, and $0.085$\% and
$0.043$\% with 2.8 m; the contours of its Fig.~1 are centred near 2.4 m.  The 2.8 m of
the text is a misprint.  A point
of a flux surface lies in a ripple well when
\begin{equation}
  \alpha \equiv \frac{\epsilon|\sin\theta|}{Nq\delta(r,\theta)} < 1,
\end{equation}
and $g(\rho)$ (\texttt{rip\_rat}) is the fraction of the poloidal circumference where
it holds, found by searching the rims on the low- and high-field sides.
$g$ is not zero near the axis: close to the low-field-side midplane
$\sin\theta\to0$ and a narrow arc always satisfies $\alpha<1$ (for the JT-60U
case below $g = 0.22$, 0.06 and 0.43 at $\rho = 0.04$, 0.25 and 0.96).  V5.61
counts these narrow core wells as the criterion gives them; no cut-off is applied.

\subsection{Ripple-trapped beam density: conservation form}

$\nbrp$ (LQr1) is the density of the ripple-trapped beam ions inside the wells;
the inventory per unit volume is $N_\r{r} = g\nbrp$.  V5.61 solves the
conservation law of $N_\r{r}$ with $\nbrp$ as the unknown,
\begin{equation}
  \ddt(g\nbrp) + \ddV\left(g U \nbrp\right) - \ddV\left(g K \ddV\nbrp\right)
  = g\frac{\Pbperp}{\Eb} + g\nueff(\sqrt\delta\,\nb - \nbrp) - g\nub\nbrp ,
  \label{eq:v561nr}
\end{equation}
\begin{equation}
  U = \aves{|\nabla V|}\ubrp, \qquad K = \aves{|\nabla V|^2}\Dbrp .
\end{equation}
$g$ stays inside every divergence.  Written for $N_\r{r}$, the diffusive flux is
$-gK\partial_V(N_\r{r}/g) = -K\partial_V N_\r{r} + K N_\r{r}\partial_V\ln g$; the
second part, the $dg/dV$ term, is what makes the flux vanish where the wells end
($g\to0$ with $\nbrp$ finite), so it cannot be dropped: without it the
ripple-trapped ions diffuse out of the well region (a test with
$-K\partial_V N_\r{r}$ alone lost 80\% of the inventory through the core-side rim).
The form (\ref{eq:v561nr}) contains no $1/g$.

Changes from V5.60, where LQr1 was $\partial_t\nbrp + \partial_V(r\ubrp\nbrp) - \partial_V(K\partial_V\nbrp) = \ldots$
without $g$:
\begin{itemize}
\item the $dg/dV$ terms (above) were missing and $\nb + g\nbrp$ was not conserved:
  the particle inventory of $\nb + g\nbrp$ was off by 2--9\% of the NBI source;
\item the advection used $r\ubrp$ (a remnant of the $\psi = r^2$ coordinate of the
  earlier code) instead of $\aves{|\nabla V|}\ubrp$, a factor $4\pi^2R$ too weak in
  circular geometry;
\item the source had no $g$, so $\nbrp$ grew where there are no wells.
\end{itemize}
Discretization: the mass, source, exchange and thermalization terms are the
element matrices $\int g\,(\ldots)\,w\,dV$; the exchange terms are the same matrices
as in LQb1 with opposite sign, so $\nb + g\nbrp$ is conserved exactly by the
discrete equations.  The advection is in strong form, $-\int\partial_V(gU\nbrp)w\,dV$;
its sum over the test functions is the wall flux $[gU\nbrp]_{V_b}$, the ripple loss
out of the plasma, and no boundary condition is needed at either end ($gU = 0$ at
the axis).  The diffusion has the natural condition.  A node with no well on
either side gets the identity $\partial_t\nbrp = 0$.  No SUPG: the element P\'eclet
number is 15--30 at the edge and $10^2$--$10^3$ inside, but the solution shows no
grid-scale structure and converges with the mesh (Sec.~\ref{sec:v561val}).

Every other equation already weighted $\nbrp$ by $g$ (LQb1, Poisson, LQi1); V5.61
also weights $Z_\r{eff}$ and the fast-ion inventory (\texttt{ANFINT}) by $g$.

\subsection{Ripple momentum closure of the beam (B-full)}

In V5.60 the parallel beam momentum (LQb3) had the two ripple terms of the 2008
model \cite{NF2008},
\begin{equation}
  -\nueff g\sqrt\delta\,\mb\nb\aves{B\upara{b}} , \qquad
  +\ddV\aves{|\nabla V|^2}\aves{B\upara{b}}\,f_t\Db\,\ddV\nb ,
\end{equation}
the momentum removed by ripple trapping and the momentum carried by the ripple
diffusion flux; the toroidal momentum (LQb4) had neither.  V5.61 removes both,
and adds nothing to LQb4.  The reason is the order of these terms in the
velocity moments of the beam distribution, not the numerical instability that
the second term caused on linear elements:
\begin{itemize}
\item The ions scattered into the ripple wells have $|\upara{}|/v < \sqrt\delta$.
  The momentum they take is $\int_{|\xi|<\sqrt\delta} m v\xi\,\nu f\,d^3v$, bounded by
  $\sqrt\delta\,mv$ per ion and $O(\delta)\,mv$ for a smooth $f(\xi)$, against
  $m\upara{b} = O(1)\,mv$ per ion in the 2008 form.  The same holds for $mRv_\zeta$
  up to the drift, $O(\rho^*_\r{b})$.
\item Ripple diffusion acts on the banana ions.  Their bounce-averaged
  distribution is even in $\xi$, so the parallel momentum flux
  $\int m v\xi\,\Gamma_r[f]\,d^3v$ vanishes at zero orbit width; the first
  non-zero part comes from the offset of the two legs of the banana,
  $O(\rho^*_\r{b})$.  For the toroidal momentum,
  $\aves{mRv_\zeta}_\r{b} = P_\zeta + Ze\aves{\psi}_\r{b}$ with $v_\para = 0$ at the tips and
  $\aves{v_\para}_\r{b} = 0$, again $O(\rho^*_\r{b})$.
\item The beam fluid of TX keeps no finite-orbit-width physics (the beam viscosity
  is zero), so keeping these $O(\delta)$ and $O(\rho^*_\r{b})$ terms alone would be
  inconsistent.  $\rho^*_\r{b} = 0.1$--0.2 at the edge of JT-60U for 80 keV D: the
  terms are dropped because the closure does not retain that order, not because
  they are negligible.
\end{itemize}
Both components vanish, so the closure also respects the radial force balance
(LQb2): a ripple term in LQb3 alone forces a spurious radial beam flux through the
multiplier $\Gamma_\r{b}$.  Ripple particle transport is kept: the trapping,
detrapping and diffusion of $\nb$ in LQb1, and LQr1.  The main effect of the ripple
in the 2008 work, the torque of the return current, comes from the charge that
these particle fluxes move and is unchanged (Sec.~\ref{sec:v561val}).

\subsection{Radial beam flux at the wall}
\label{sec:beamwall}

With \texttt{MDBEAM} = 1, $\Gamma_\r{b} = \nb\uV{b}$ (LQb2) is the Lagrange multiplier of
the beam radial force balance; V5.60 had no condition on it at $\rho_b$.  V5.61
imposes
\begin{equation}
  \Gamma_\r{b}(\rho_b) = 0 .
\end{equation}
This is a closure of the reduced model, not a statement that fast ions are not
lost to the wall.  The model has no surface or orbit-loss channel for the beam;
outside the separatrix the beam is lost by the volume loss $\nubL\nb$.  The wall
value of the multiplier was not an outflow law: it was whatever the force balance
at the wall node required, with the diamagnetic term of a beam with fixed
$T_\r{b} = 2\Eb/3$, whereas the fast-ion orbit width (10--20 cm) exceeds the gradient
lengths of the SOL.  Left free, that value pushed $\nb$ at the wall node towards zero
and made the Picard iteration cycle without converging (a period-5 cycle of
$\Gamma_\r{b}$ on the circular test equilibrium).  The alternatives were rejected: the
SOL beam removal used with \texttt{MDBEAM} = 0, 2 rewrites $\nb$ outside the equations
and contradicts the constraint rows; an outflow $\Gamma_\r{b} = \nb v_\r{out}$ needs a
physical $v_\r{out}$ and would double-count $\nubL$.  The condition changes runs with
NBI by 1.5--5\% in the core rotation for perpendicular injection and by 0.5\% for the
JT-60U case below.  It should be revisited if a finite-orbit-width or orbit-loss
model is introduced.

\subsection{Beam--thermal friction}
\label{sec:v561friction}

The friction between the beam and the thermal species follows the fast-ion
coefficients of Kikuchi and Azumi \cite{KA1995} (\texttt{imodel\_neo(3)} = 2, isotropic fast
ions in the pitch-angle-scattering regime): $\ell_{11}^{\r{fe}} = \ell_{11}^{\r{ef}}$,
$\ell_{11}^{\r{f}k} = \ell_{11}^{k\r{f}}$ ($k$ = i, z),
$\ell_{11}^{\r{ff}} = -\sum_k\ell_{11}^{\r{f}k}$, the Maxwellian $\ell_{11}^{kk}$ reduced
by $\ell_{11}^{k\r{f}}$, and the heat friction $\ell_{21}^{\r{ef}} = -\frac32\ell_{11}^{\r{ef}}$ with
$\ell_{21}^{\r{fe}} = 0$.  The 11 block is symmetric, so the friction forces sum to zero;
only the heat friction is not symmetric, as the paper states.  In V5.60 the
fast-ion row was obtained from the thermal row with the Braginskii ratio
$\tau_{ee}/\tau_{ii}$, which assumes equal Coulomb logarithms, and converted back to
physical units with the actual $\tau_{ee}$ and $\tau_{ii}$.  The result was
$\ell_{11}^{\r{fi}}/\ell_{11}^{\r{if}} = \ln\Lambda_{ee}/\ln\Lambda_{ii}$ (0.82--0.86 for D,
1.00--1.08 for C) and a net force of 2--10\% of the fast-ion friction: 3\% of the
NBI torque in the JT-60U case, 13\% at lower $T_\r{e}$.  V5.61 converts with the same
collision times; the column sums of the 11 block are now $10^{-16}$ of the diagonal.
Applicability: the coefficients assume an isotropic fast-ion distribution and the
slowing-down regime $v_\r{Ti} \ll v_\r{b}$, and they ignore the fast-ion heat flow;
for a tangential beam the distribution is anisotropic.

\subsection{Discrete momentum conservation of the friction}
\label{sec:frictionfem}

The friction force on species $a$ is $\sum_b \ell^{ab}\,u_b$ with $\sum_a\ell^{ab} = 0$,
so the forces cancel for any flows.  The coefficients satisfy this at every node (to
$10^{-8}$ of the diagonal, MI).  The equations of species $a$ are written per unit
mass density, so the term of row $a$ acting on the momentum $X_b = n_b\aves{Ru_b}$ has
the coefficient $\ell^{ab}/(m_a n_a)\times n_a/n_b$; up to V5.61.00 it was integrated
as the product of the two interpolants (\texttt{fem\_int(20,lab,Ratden)}, and
\texttt{fem\_int(20,lab,n)} in the parallel equations), while the diagonal term of
the same column used the interpolated coefficient (\texttt{fem\_int(2,lab)}).  Per
element the two differ by $-h\,\Delta a\,\Delta b\,(X_1+X_2)/12$, where $\Delta a$,
$\Delta b$ are the changes of the two factors across the element, so the columns did
not sum to zero and the friction left a net force of $O(h^2)$.  It was measured
directly: the residual falls by 2.2--2.3 from \texttt{NRMAX} = 60 to 90 and by 3.8--3.9
to 120, and in the JT-60U case it came almost entirely from the ion--impurity pair
near $\rho = 0.8$, where $n_i/n_z$ and the coefficients vary most.  Every friction
term of a column is now assembled from the nodal product of its row-dependent
factors ($\ell^{ab}n_a/n_b$, $\ell^{ab}n_a$, and $\ell_{12}^{ab}n_a$ for the heat
friction, with the row-independent factors $I/\aves{B^2}$ and $\partial_V T_b$
integrated alike in every row).  The element matrices then sum to zero over a column
whenever the nodal coefficients do, and the continuous equations are unchanged.  The
change of the solution is $O(h^2)$: in the JT-60U case at 1.5 s, $u_{\r{i}\phi}(0)$
goes from 658 to 656 km/s, the total angular momentum by $-0.2$\%, the impurity density
by up to 1\% and $E_r$ by up to 0.7\% (normalised to their maxima); without NBI, where
the spurious force was comparable to the whole toroidal angular momentum, the small
rotation changes by up to 5\%.  The Picard counts are unchanged.

\subsection{Validation}
\label{sec:v561val}

\paragraph{Regression.}  Without NBI all runs are bit-identical to V5.60 (see the
start of this section).  With NBI and \texttt{FSRP} = 0 the solution changes by the
intended $\Gamma_\r{b}(\rho_b)=0$ and the friction correction: on the X-point test
equilibrium with 5 MW perpendicular NBI the total toroidal angular momentum at
0.5 s is 16\% lower (the injected torque is zero there, so the spurious friction
source was comparatively large), with 3 MW co-tangential plus 1 MW perpendicular
2\% lower; bulk densities and temperatures change by less than $10^{-3}$.  On the
circular test equilibrium the wall cycle of $\Gamma_\r{b}$ is gone (V5.60: 25 of 5000
steps at \texttt{ICMAX}; V5.61: none).

\paragraph{Stability matrix.}  \texttt{FSRP} = 1, perpendicular NBI, \texttt{NANDER} = $-20$,
X-point and circular, 0.3 s (0.15 s for $\Delta t = 2.5\times10^{-5}$):
\begin{center}
\begin{tabular}{|l|c|c|}\hline
case & X-point: steps at \texttt{ICMAX} / mean & circular \\\hline
NRMAX 60, $\Delta t = 10^{-4}$, 2.5 / 5 / 10 MW & 1 / 15,\ 0 / 21,\ 0 / 28 & 0 / 19,\ 0 / 28,\ 0 / 34 \\
$\Delta t = 2.5\times10^{-5}$, $2\times10^{-4}$ & 0 / 11,\ 0 / 28 & 0 / 14,\ 1 / 37 \\
NRMAX 120 & 13 / 42 & 1020 / 68 \\\hline
\end{tabular}
\end{center}
No run stops and none shows a grid-scale mode: the odd-even part of the beam
parallel momentum near the edge is 0.016--0.026 of its maximum at \texttt{NRMAX} = 60,
the level of the non-uniform mesh, and 0.005 at 120.  The steps at \texttt{ICMAX} on
the circular equilibrium at \texttt{NRMAX} = 120 are limited by the impurity
diamagnetic flow (LQz8) at the axis node.  They do not occur with \texttt{FSRP} = 0 and
remain when the ripple terms inside $\rho = 0.3$ are removed (tested: $g = 0$ or
$\Dbrp = 0$ there): it is the known weakness of the impurity equations at the axis,
triggered by the state the ripple produces, not a defect of the ripple model.

\paragraph{Conservation.}  Volume integrals of every term of every equation
(JT-60U case below, every step): the discrete time derivative equals the sum of
the terms to $10^{-6}$ for all density and momentum equations; $\nb+g\nbrp$ closes
with a residual below $10^{-5}$ of the NBI source.  The injected torque is 6.73 N\,m
(analytic $2P R_\r{tan}/v_\r{b}$: 6.76 N\,m).  The friction between the beam and the
thermal species conserves momentum since the correction of $\ell_{11}^{\r{f}k}$
(Sec.~\ref{sec:v561friction}), and the friction terms of all species sum to less
than $10^{-8}$ N\,m at all times since the assembly of Sec.~\ref{sec:frictionfem}; before
that, the friction between the thermal ions and the impurity left a discrete
residual that grew with the rotation to $+0.037$ N\,m (0.55\% of the injected
torque) at 1.5 s.  The viscous flux at the wall is zero by the weak form; the rotation
is damped by the SOL parallel loss and charge exchange at the edge.

\paragraph{Ripple benchmark.}  JT-60U-like circular case of \cite{NF2008}: $R=3.4$ m,
$a=0.95$ m, $B=2.6$ T, $I_\r{p}=1$ MA, $\kappa=1.15$ for the ripple, $N=18$,
$\delta_\r{min}=2\times10^{-6}$ (edge ripple 1.26\%), 3.2 MW co-tangential and 0.4 MW
perpendicular NBI, $\chi = \mu = 0.22(1+19\rho^2)$ m$^2$/s, $D = 0.15(1+4\rho^2)$ m$^2$/s,
1.5 s:
\begin{center}
\begin{tabular}{|l|c|c|c|c|c|c|}\hline
case & $u_{\r{i}\phi}(0)$ & $u_{\r{i}\phi}(0.82)$ & $E_r(0.5)$ & $j_r(0.97)$ & $j\times B$ torque & $dL_\phi/dt$ \\
 & km/s & km/s & kV/m & A/m$^2$ & N\,m & N\,m \\\hline
no ripple & 658 & 152 & $+83$ & $-0.020$ & $+0.18$ & $+0.12$ \\
half ripple (FST) & 417 & 60 & $+34$ & $-0.108$ & $-3.40$ & $+0.05$ \\
full ripple & 128 & $-42$ & $-22$ & $-0.185$ & $-7.44$ & $0.00$ \\\hline
\end{tabular}
\end{center}
An earlier version of this table gave the state at 0.6 s.  At that time the runs
without ripple and with half the ripple were still transient ($dL_\phi/dt = +1.57$
and $+0.87$ N\,m against the injected 6.73 N\,m; $u_{\r{i}\phi}(0)$ = 473 and
325 km/s); the table now gives the quasi-steady state at 1.5 s, where
$dL_\phi/dt$ is below 2\% of the injected torque in all three runs.  With ripple
the return current and its $j\times B$ torque changed by 5--11\% from 0.6 s to 1.5 s,
so the ripple results are qualitatively unchanged; the co-rotation without ripple
is 40\% higher than at 0.6 s.
The return current grows with the ripple and has the magnitude of \cite{NF2008}
($-0.08$ and $-0.16$ A/m$^2$ with and without FST), its torque is counter and
grows, $E_r$ moves negative and the edge rotates counter with the full ripple;
halving the ripple weakens all of these.  The qualitative physics of the 2008
work is kept by B-full and the conservative LQr1.  The co-rotation without ripple
is 2.2 times that of \cite{NF2008} (300 km/s).  The injected torque agrees with
its analytic value, so this is a question of the momentum damping and transport
(the losses concentrate at the edge in the V5.60 SOL and neutral models), not of
the ripple model.



\section{Models, switches and defaults of V5.60}
\setcounter{equation}{0}

V5.60 changes the default SOL and neutral physics.  The detailed physics,
the verification and the development history are in
\texttt{neutral\_transport.tex}; this section states what the code solves.

\subsection{SOL parallel loss}

Outside the separatrix the parallel loss rate is
\begin{equation}
  \nuL = k_L \frac{C_s}{L_c}\, g(\rho),\qquad L_c = 2\pi q R,\qquad
  g = \frac{x^2}{1+x^2},\quad x = \frac{r-a}{\delta}\ (r\ge a),\ x=0\ (r<a),
\end{equation}
and enters every SOL equation in Robin form, e.g.\ $-\nuL(n_s-\nsup{s})$.
$k_L$ is \texttt{FSLP}.  The ramp $g(\rho)$ is numerical smoothing: it keeps
$\nuL$ from jumping at $\rho=1$, which the FEM cannot represent, and its
continuum limit is the step.  \texttt{MDLSOL} selects its width $\delta$:

\begin{center}
\begin{tabular}{|c|l|}\hline
\texttt{MDLSOL} & $\delta$ \\\hline
0 & $3\sqrt{0.5\,T_i(0)\,m_i}\,q(a)/(Z_ieB_\phi(a))/\sqrt{a/R}$ (V5.52) \\
1 & ion banana width at the separatrix, $\sqrt{\epsilon}\,\sqrt{T_i m_i}/(Z_i e B_\theta)$ (default) \\\hline
\end{tabular}
\end{center}

With \texttt{MDLSOL} = 0 the width is evaluated with half the on-axis ion
temperature and the edge $q$; for the X-point test equilibrium distributed with
the code (\texttt{eqdata}, $R=3.2$ m, $a=0.8$ m) it is 1.8 times the whole SOL, so
$g$ is only 0.17 at the wall and the ramp suppresses the loss everywhere instead
of smoothing its onset.  \texttt{MDLSOL} = 1 evaluates the banana width with the
local ion temperature and poloidal field at the separatrix and gives
$g\simeq 0.99$ at the wall.

The default $k_L = 1$ is the value of the sheath-limited two-point model
(pressure conservation, Mach number 1 at the target, upstream and target
densities related by $n_\r{t}=n_\r{u}/2$) when $2\pi qR$ is read as the length
between the two targets: a flux tube of that length loses $2n_\r{t}C_s$ per unit
area while holding $n_\r{u}\,2\pi qR$, so $\nuL = C_s/(2\pi qR)$.  V5.52 used the
empirical value $k_L=0.3$, a typical upstream Mach number of the SOL flow; with
\texttt{MDLSOL} = 1 it does not give a stable SOL on the X-point test
equilibrium.

\subsection{Wall boundary condition for the densities}
\label{sec:wallbc}

The table of the previous section lists a Neumann condition for the densities
at the wall $\rho=\rho_b$, said to come naturally from the pressure equation.
The natural boundary condition of that equation is, however, zero heat flux, not
a condition on the density.  Up to V5.52 the densities therefore had no
condition at the wall, and the radial particle flux through $\rho_b$ was fixed by
the last element.  On the X-point test equilibrium the ratio of the ion flux
through the wall to that through the separatrix,
$\Gamma_\r{wall}/\Gamma_\r{sep}$, was $+0.57$, $-0.40$ and $-0.28$ at
\texttt{NRMAX} = 60, 120 and 480: depending on the mesh the wall acted as a
source or a sink of particles, and refining the mesh did not make it vanish.  The radial flux
$n_s\uV{s}$ cannot take the condition instead: its equation is first order and
not integrated by parts, so it already has its one condition at the axis.
For the charged species $\rho_b$ is a computational boundary at the outer edge of
the SOL, not a material wall: their exhaust is carried by $\nuL$.

V5.60 (\texttt{MDLWBC} = 1) imposes $dn_s/d\rho = 0$ on $\nee$, $\ni$ and $\nz$.
In the P1 basis this is the constraint
\begin{equation}
  n_s(N) = n_s(N-1),\qquad N = \texttt{NRMAX}.
\end{equation}

\paragraph{Constrained degrees of freedom, in general.}
For a linear system $A x = b$ with the constraint $x_N = x_{N-1}$, write
$x = T y$, where $T$ is the identity on $y_1\dots y_{N-1}$ and copies $y_{N-1}$
into $x_N$.  The constrained Galerkin system is
\begin{equation}
  T^\r{T} A T\, y = T^\r{T} b .
\end{equation}
$T^\r{T}$ adds row $N$ to row $N-1$; $AT$ adds column $N$ to column $N-1$.  In
finite elements this is the basis reduction $\phi_{N-1}\to\phi_{N-1}+\phi_N$:
the trial and test functions are constant on the last element, which is the
P1 form of $dn/d\rho=0$.  The code keeps the system square and solves the
equivalent
\begin{equation}
  \mbox{row } N-1 \leftarrow \mbox{row } N-1 + \mbox{row } N,\qquad
  \mbox{row } N \leftarrow x_N - x_{N-1} = 0 ,
\end{equation}
with $x_N$ still an unknown, so the other equations keep their coupling to it and
the column sum is carried by the constraint row.
Since $\sum_i\phi_i = 1$, the sum of all rows of a continuity equation is the
global particle balance; each column of $T^\r{T}$ sums to one, so that sum,
and with it particle conservation, is unchanged.  Replacing row $N$ by the
constraint without adding it to row $N-1$ would drop the particle balance of the
wall half-cell, an $O(h)$ conservation error.

\paragraph{Implementation.}
\texttt{txcala.f90}, subroutine \texttt{BOUNDARY\_ZEROGRAD(LQ)}, called from the
boundary-condition block of \texttt{TXCALA} for \texttt{LQe1}, \texttt{LQi1} and
\texttt{LQz1} after the Dirichlet eliminations.  In the band storage the
coefficient of term \texttt{NC} of equation \texttt{NQ} in row \texttt{NR} is
\texttt{ALC}, \texttt{BLC} or \texttt{CLC}\texttt{(NR,NC,NQ)} for
\texttt{X(NR+1)}, \texttt{X(NR)} or \texttt{X(NR-1)}, and \texttt{PLC} holds the
known terms.  For every term \texttt{NC = 0..NLCMAX}:
\texttt{BLC(N-1)~+=~CLC(N)}, \texttt{ALC(N-1)~+=~BLC(N)},
\texttt{PLC(N-1)~+=~PLC(N)}; row \texttt{N} is then cleared and set to
\texttt{BLC(N,1) = 1}, \texttt{CLC(N,1) = -1}.  The time term and its
right-hand side are built from the same coefficient arrays, so they are folded
with them.  The condition converges at first order in the mesh size:
$\Gamma_\r{wall}/\Gamma_\r{sep}$ is $-1.8$, $-1.1$, $-0.5$ and $-0.3$\% at
\texttt{NRMAX} = 60, 120, 240 and 480.

\subsection{C$^0$--C$^{5+}$ superstage with a source-slaved C$^0$ part}

Carbon is carried by two densities: $\nz$ for C$^{6+}$ and the superstage
$\nnz = \sum_{q=0}^{5} n_{\r{C}^{q+}}$, called the bundle below.  The bundle
exchanges particles with C$^{6+}$ through the effective ionization $f_5S_5$
and recombination $A_6$ of Sec.~1, and is fed by sputtering in the SOL:
$\gamma_\r{0z}$ (\texttt{rGamm0zz}) is the fraction of the C$^{6+}$ parallel loss
and $\gamma_\r{i,z}$ (\texttt{rGamm0iz}) the fraction of the main-ion parallel
loss that return as sputtered carbon.

In local ionization balance the C$^0$ fraction of the bundle is below
$2\times10^{-6}$, so the bundle is almost entirely magnetized C$^{1+}$--C$^{5+}$
ions.  With \texttt{MDLNZD} = 1 it is therefore diffused with the anomalous
particle diffusivity of the charged species, $D_{n\r{z}} = D_\r{c}$ (\texttt{Dfs}).
The exception is the carbon that has just been sputtered, which enters as
neutral C$^0$.  Its lifetime $1/(S_0\nee)$ is of the order of microseconds,
shorter than every other time scale here, so its density follows its local
source balance:
\begin{equation}
  n_\r{C0} = \min\left(\frac{S_\r{wall}}{S_0\nee},\ \nnz\right),\quad
  S_\r{wall} = \gamma_\r{0z}\nuL \max(\nz-\nsup{z},0) + \gamma_\r{i,z}\nuL\max(\ni-\nsup{i},0),
\end{equation}
where $S_0$ is the C$^0\to$C$^{1+}$ ionization rate (OpenADAS scd96,
\texttt{data/Carbon-S0F0\_20260921.txt}).  $n_\r{C0}$ enters the $\nnz$
equation as a known field through the excess flux with
$\Delta D_\r{C0} = D_\r{C0} - D_\r{c}$, the part of the C$^0$ flux
$-D_\r{C0}\nabla n_\r{C0}$ that the $D_\r{c}$ term does not already carry, so that the
bundle flux is exactly $-D_\r{c}\nabla(\nnz-n_\r{C0}) - D_\r{C0}\nabla n_\r{C0}$.
Here $D_\r{C0}$ is the Larsen-limited (see below) value of $\bar v_{0z}^2/(3S_0\nee)$,
the short-mean-free-path diffusivity of C$^0$, and $\bar v_{0z} = \sqrt{4V_{0z}^2/\pi}$
its mean speed ($V_{0z}$ = \texttt{V0z}), and through the parallel loss of the
charged part of the bundle, $(1-\gamma_\r{0z})\nuL(\nnz-n_\r{C0})$, of which the
fraction $\gamma_\r{0z}$ is sputtered back (the $\nnz$ equations of Secs.~1 and 3).
$n_\r{C0}$ is evaluated from the previous Picard iterate, $S_\r{wall}$, $S_0\nee$ and the
bound $\nnz$ alike, so the bound $n_\r{C0}\le\nnz$ applies to the bundle at the time level
being solved for, and the $\nnz$ equation stays linear in $\nnz$.  With \texttt{MDLNZD} = 0 (the V5.52 model),
$D_{n\r{z}} = D_{01}$, the diffusivity of the slow hydrogen neutrals, and the
$n_\r{C0}$ terms are absent.

\paragraph{Element-based C$^0$ limiter.}
$D_\r{C0}$ is Larsen-limited, $D\to D/(1+D|\nabla n_\r{C0}|/(\alpha\bar v_{0z}n_\r{C0}))$,
$\alpha$ = \texttt{FLlim}.  The limiter is evaluated on each element from the
element mean of $n_\r{C0}$ and the element gradient.  It acts on the whole
$D_\r{C0}$, because it bounds the whole C$^0$ flux by free streaming, and
$D_\r{c}$ is subtracted afterwards with its sign: $\Delta D_e = D_e - D_{\r{c},e}$, where
$D_e$ is the limited $D_\r{C0}$ and $D_{\r{c},e}$ the element mean of $D_\r{c}$.
$\Delta D_e<0$ is not a negative diffusion: it removes the part $D_{\r{c},e}$ of the C$^0$
flux that the $D_\r{c}$ term already carries, and since $n_\r{C0}$ is a known field
the matrix is unchanged.  It occurs in the first element outside the separatrix,
where $n_\r{C0}$ rises from zero within one element and the limited $D_\r{C0}$ falls to
about $\alpha\bar v_{0z}h/2$ ($h$ the radial width of the element): at NRMAX = 480,
not at NRMAX $\le$ 240.  Clipping $\Delta D_e$ at zero there would let the C$^0$ flux
exceed free streaming by 7\%.  The element flux
$F_e = \aves{|\nabla V|^2}_e \Delta D_e\,\Delta n_\r{C0}/\Delta V_e$ enters the two
rows of the element as $+F_e$ and $-F_e$, so the term is exactly conservative.
The element evaluation is needed because inside the separatrix the source, and
hence $n_\r{C0}$, vanishes: a nodal ratio $|\nabla n_\r{C0}|/n_\r{C0}$ would be
$0/0$ at the separatrix node, and whether that node is limited would be left to
rounding.  The limiter is part of the model and is always applied with
\texttt{MDLNZD} = 1.  Without it $D_\r{C0}$ becomes very large in the SOL, and the
divergence of the known C$^0$ flux, a source that does not scale with $\nnz$,
drives $\nnz$ negative just outside the separatrix.

\paragraph{Validity of the source-slaved C$^0$ part.}
The slaving of C$^0$ to its source is implemented and is part of the production
model.  What remains open is the range of validity of this quasi-steady
approximation and the lack of a formal positivity guarantee.  The mean free path
of C$^0$ is comparable to the width of the C$^0$ layer, so the local source
balance does not resolve the inner structure of the layer.  Positivity of $\nnz$
is kept by the limiter in the cases computed, but it is not guaranteed, because
the divergence of the known C$^0$ flux is a source that does not scale with $\nnz$.
The bound $n_\r{C0}=\nnz$ marks the limit of the model.  It is never reached in
the converged iterates of the production cases; it is reached only in the first
sweep of a run started from the initial profiles, where $\nnz$ is still close to
zero.  Where it holds persistently, the lagged treatment of $n_\r{C0}$ is not
appropriate: the C$^0$ flux in the bounded nodes is then a diffusion of $\nnz$ itself
with $\Delta D_\r{C0}\approx 7$--$15\,D_\r{c}$, taken from the previous iterate, and the
iteration diverges (forced with a C$^0$ ionization rate reduced 100-fold; with the
bound taken from the previous time step, the time stepping diverges instead).
\texttt{TXLOOP} warns when the bound is still active in the last Picard sweep of a
step, and with \texttt{IDIAG} $\ge$ 2 prints the number of bounded nodes at every
step.  Such conditions call for an independent C$^0$ transport equation, an
implicit treatment of the C$^0$ flux, or a kinetic (Monte-Carlo) treatment, which
would also delimit the validity of the approximation.

\subsection{Larsen limiter on the hydrogen slow neutrals (optional)}

The slow hydrogen neutrals are diffused with the short-mean-free-path
diffusivity $D_{01} = \bar v_0^2/(3S_\r{tot}) = \lambda\bar v_0/3$, $S_\r{tot}$ being the
ionization plus charge-exchange rate and $\lambda = \bar v_0/S_\r{tot}$ the mean free
path.  The Larsen limiter replaces $D$ by $D/(1+R)$ with
\[
  R = \frac{D|\nabla n_{01}|}{\alpha\,\bar v_0\, n_{01}}
    = \frac{\mbox{flux the diffusion law demands}}{\mbox{free-streaming flux it may not exceed}}
    = \frac{\lambda}{3\alpha L_n},\qquad \alpha = \texttt{FLlim},
\]
$L_n$ being the gradient length: the limited flux $D|\nabla n|/(1+R)$ is the
diffusion flux where $R\ll1$ and $\alpha\bar v_0 n$ where $R\gg1$.  $\alpha = 1/4$ is
the one-sided flux of an isotropic population: for particles of speed $\bar v$
spread uniformly over all directions, those crossing a plane in one direction
carry $n\bar v\int_0^1\mu\,d\mu/2 = n\bar v/4$.  $\alpha = 1/2$ would correspond to
all particles moving into one half-space, uniformly in angle, and $\alpha = 1$ to
all of them moving in the same direction, the kinematic upper bound.  In the
V5.60 solution without the limiter $R$ is 1--3 in the SOL, about 2 at the
separatrix and up to 6 at $\rho\simeq0.8$--0.9 ($\mathrm{Kn} = \lambda/L_n\sim1$): the
diffusion approximation is marginal wherever the neutrals matter.

\texttt{MDLN0FL} = 0 (default) leaves $D_{01}$ unlimited.  \texttt{MDLN0FL} = 2 evaluates
the limiter element by element, as for the C$^0$ flux, from the element mean
$\bar n$ and the element difference $\Delta n$ of the current iterate; since
$|\Delta n|/\bar n\le2$ on an element, the factor is bounded by
$1 + 2\lambda/(3\alpha h)$.  \texttt{MDLN0FL} = 1 is the earlier nodal form, with
$|\nabla n|/n$ from the spline derivative, kept only to reproduce earlier runs.
The limited ionization front at $\rho\simeq0.85$--0.9 is exponential with the decay
length $\alpha\lambda\simeq5$\,mm, shorter than an element at \texttt{NRMAX} = 60
(about 1\,cm).  On that unresolved front, and in the tiny-$n_{01}$ tail inside it,
the nodal ratio is large and erratic, and the Picard iteration does not settle:
with \texttt{MDLN0FL} = 1, 496 of 500 steps reach \texttt{ICMAX} at \texttt{NRMAX} = 60
(\texttt{EPS} = $10^{-6}$), the neutral-only relaxation \texttt{RLXN0} = 0.5 (used only with
\texttt{MDLN0FL} = 1) leaves 2--3.  \texttt{MDLN0FL} = 2 needs no relaxation: 20.0 / 21.7
sweeps on average and no step at \texttt{ICMAX} (X-point / circular, \texttt{NRMAX} = 60),
the same as without the limiter.  The two forms agree to 0.5--2\% where both
converge.  Against the limiter off, the limiter raises the neutral inventory by
about 30\% and changes $\nee$ at the separatrix and at the wall by $-3$\% and
$+4$\%; $Z_\r{eff}(0)$ is unchanged.

The limiter is a bound on the diffusion flux, not a substitute for kinetic
neutral transport.  A straight-flight reference (the same sources, attenuated at
$S_\r{tot}$) ionizes or charge-exchanges 35\% of the slow neutrals inside the
separatrix (X-point, \texttt{NRMAX} = 240); plain diffusion gives 49\% and the limiter
with $\alpha = 1/4$ 25\%, with almost nothing beyond $\rho = 0.95$ against 12\% in the
reference.  $\alpha = 1/2$ recovers the total but not the ballistic tail, and no
single $\alpha$ recovers both.  The difference between limiter on and off is of the
size of the uncertainty from the neutral energy alone (the reference gives 24\%
and 48\% at 0.5 and 4.5\,eV).  The limiter is therefore off by default.

\subsection{Anderson acceleration of the Picard iteration (optional)}
\label{sec:anderson}

Each time step is solved by Picard iteration: the coefficients are evaluated
from the previous iterate $x_k$ and the linear system gives $g_k = G(x_k)$.  On
fine meshes the iteration has a slow mode localised on the first few nodes at
the magnetic axis, in $BV_{1s}$ (\texttt{LQ*8}) of all species, whose contraction
factor per sweep tends to 1 as the axis element shrinks (0.70, 0.91, 0.98 at
\texttt{NRMAX} = 120, 240, 480 with $\Delta t = 10^{-4}$); at \texttt{NRMAX} = 480
it takes about 500 sweeps to reach \texttt{EPS} = $10^{-6}$.

With \texttt{NANDER} $\ne 0$ the next iterate is instead
\[
  x_{k+1} = g_k - \sum_j \gamma_j\,(\Delta x_j + \Delta f_j),\qquad
  \gamma = \arg\min\Bigl|f_k - \sum_j \gamma_j \Delta f_j\Bigr|_2,
\]
with $f = g - x$ and $\Delta x_j$, $\Delta f_j$ the differences over the last
$|\texttt{NANDER}|$ sweeps of the step (Anderson mixing with $\beta = 1$).  Each
variable is scaled by its profile RMS at the first sweep of the step, as in the
convergence test.  With \texttt{NANDER} $< 0$ the rows of \texttt{LQm1}--\texttt{LQm5} do not
enter the least-squares fit, because near convergence their increments are
rounding (for $\psitd$ the floor of Sec.~\ref{sec:lqm2tol}); the correction is
applied to all variables.  The history
is restarted at every time step, and dropped (the plain $g_k$ is used) if the
least-squares solve fails or returns $|\gamma| > 10^3$.  Once the adaptive
relaxation of the step has found the iteration not contracting, the rest of the
step is left to the relaxation alone: the linear model behind the mixing does not
hold there, and with the mixing kept on the first steps from the initial profiles
at \texttt{NRMAX} = 960 end in a negative density.  A fixed point of $G$ is
a fixed point of the accelerated iteration, so the converged solution is that of
plain Picard; \texttt{NANDER} = 0 (default) is bit-identical to the plain iteration.

With \texttt{NANDER} = $-20$ (100 steps from $t = 0.02$, $\Delta t = 10^{-4}$,
\texttt{EPS} = $10^{-6}$, \texttt{ICMAX} = 100, X-point equilibrium) the average
number of sweeps falls from 22.7, 36.4, 99.4 to 16.0, 28.4, 58.8 at \texttt{NRMAX} =
60, 120, 240, and the steps that reach \texttt{ICMAX} at \texttt{NRMAX} = 240 from 86 to 9;
the circular equilibrium behaves alike.  The solution agrees with the plain
iteration to within the convergence tolerance.  At \texttt{NRMAX} = 480 the
\texttt{LQ*8} mode is removed as well.  These figures were measured with the
purely relative test of $\psitd$ (\texttt{LQm2}), under which almost every step at
\texttt{NRMAX} = 480 still reached \texttt{ICMAX} on the rounding floor of $\psitd$,
which the acceleration does not affect.  With the floor term of
Sec.~\ref{sec:lqm2tol} the steps at \texttt{NRMAX} = 240 and 480 end when the
other variables have converged: with \texttt{NANDER} = $-20$ the average number
of sweeps is 33.5 (240) and 67.1 (480), and 3 of 100 steps at 480 reach
\texttt{ICMAX}, on the other variables.

\subsection{Convergence test of $\psitd$ (\texttt{LQm2})}
\label{sec:lqm2tol}

The sweeps of a time step stop when, for every variable $X_q$,
$\max_j|\Delta X_{q,j}| \le \texttt{EPS}\,\mathrm{rms}(X_q)$, $\Delta X_q$ being the
change over the sweep.  For $X_2 = \psitd$ this relative test has a floor that no
iteration removes, and the test is
\[
  \max_j|\Delta X_{2,j}| \le \texttt{EPS}\,\mathrm{rms}(X_2)
  + C_\r{r}\,\epsilon_\r{m}\,\texttt{NRMAX}^2\,
    \frac{\mu_0\,\mathrm{rms}(X_5)}{\texttt{ADV}\,\Delta t},
  \qquad C_\r{r} = 1,
\]
with $X_5 = \psit/\mu_0$ (\texttt{LQm5}), $\epsilon_\r{m}$ the machine epsilon and
\texttt{ADV} the weight of the implicit terms (2/3 for BDF2, 1 in the first step).
\texttt{LQm2} is not removed from the test: the relative term is kept as it is
for every variable, and only the floor term is added.  The test of every other
variable (including $\psid$, \texttt{LQm3}), the equations and the fixed point are
unchanged.

The floor has two factors.  $\psit$ is mostly vacuum flux, and in Amp\`ere's
law (the \texttt{LQm2} row) the $\nabla\times B_\phi$ term applied to it is a
cancellation of terms larger than the result, $|KX_5|/(|K||X_5|) = 3.4\times10^{-7}$
at \texttt{NRMAX} = 60 falling as $\texttt{NRMAX}^{-2}$.  A rounding of relative
size $\epsilon_\r{m}$ in that row, in the assembly or in the solve, leaves $\psit$
determined only to $(0.05$--$0.07)\,\epsilon_\r{m}\texttt{NRMAX}^2$.  Faraday's
law (the \texttt{LQm5} row) then gives $\psitd$ as the difference quotient
$\mu_0(X_5 - X_5^\r{old})/(\texttt{ADV}\,\Delta t)$, which multiplies that rounding
by $\psit/(\texttt{ADV}\,\Delta t\,|\psitd|)$, $7\times10^6$ at $\Delta t = 10^{-4}$,
$t = 0.02$\,s, and without bound as $\psit$ approaches a steady state.  The
relative floor of $\psitd$ is therefore about $4\times10^{-7}$, $1.5\times10^{-6}$,
$4.6\times10^{-6}$, $1.8\times10^{-5}$ at \texttt{NRMAX} = 60, 120, 240, 480
($\Delta t = 10^{-4}$, $t = 0.02$\,s), grows as $1/\Delta t$, and grows in time:
at \texttt{NRMAX} = 60 it reaches $10^{-5}$ by $t = 0.5$\,s.  At the floor the
equations are satisfied to rounding (backward error of the \texttt{LQm2} and
\texttt{LQm5} rows $10^{-15}$--$10^{-11}$, normwise residual of every equation below
$4\times10^{-9}$); the increment of $\psitd$ is the forward error of an
ill-conditioned component, not an unconverged iteration.  The second term of the
test is that floor with a safety factor.

$C_\r{r}$ is not a physical tolerance on $\psitd$.  $C_\r{r} = 0.3$ lies on the
largest floor observed (the rounding of $\psit$ reaches 4--6.5 times its geometric
mean) and leaves no margin for rounding that depends on the compiler, the BLAS or
the order of operations; $C_\r{r} = 1$, three to four times it, is used.  The two
values give the same numbers of sweeps to within a few per cent and solutions
that cannot be told apart.  The test on $\psitd$ keeps its role while the
iteration converges, since $\psitd$ responds to a change of the currents about
$10^4$ times more strongly than they change (\texttt{NRMAX} = 60).  In the printed and returned
convergence measures the value for \texttt{LQm2} is the part of the relative
increment above the floor term.

Against a reference solved with the relative test and up to 300 (100 steps) or
200 (3000 steps) sweeps per step, \texttt{NANDER} = $-20$, the test with the floor
term gives (\texttt{EPS} = $10^{-6}$, \texttt{ICMAX} = 100, $\Delta t = 10^{-4}$)

\begin{center}
\begin{tabular}{|l|c|c|}\hline
 & relative test & with the floor term \\
 & sweeps / steps at \texttt{ICMAX} & sweeps / steps at \texttt{ICMAX} \\\hline
X-point, 240, 100 steps       & 58.8 / 9    & 33.5 / 0 \\
circular, 240, 100 steps      & 69.8 / 10   & 37.1 / 0 \\
X-point, 480, 100 steps       & 99.6 / 97   & 67.1 / 3 \\
circular, 480, 100 steps      & 99.3 / 95   & 77.0 / 5 \\
X-point, 240, 3000 steps      & 89.0 / 2167 & 29.9 / 0 \\
circular, 240, 3000 steps     & 89.6 / 2093 & 33.6 / 0 \\\hline
\end{tabular}
\end{center}

The remaining steps at \texttt{ICMAX} are steps on which another variable has not
converged.  Densities, temperatures, current and $Z_\r{eff}$ differ from the
reference by less than $7\times10^{-8}$; $\psitd$ and the poloidal electric field derived from it by
$0.2$--$1.5\times10^{-4}$, within the $0.3$--$2.1\times10^{-4}$ by which the
relative test with \texttt{ICMAX} = 100 differs from the same reference, that is,
within the floor.  The loop voltages and the Poynting-flux terms follow within
$6\times10^{-7}$ over 3000 steps, except the rate of change of the poloidal
electric-field energy, which is built on $\epsilon_0\psitd^2$ and carries the floor
of $\psitd$ under either test.  Volume-averaged densities, stored energy and
current agree to $10^{-10}$ and the toroidal angular momentum to $2\times10^{-8}$.
The BDF2 errors and orders of $n$, $T$, $j$ and $Z_\r{eff}$ are unchanged
(\texttt{NRMAX} = 60, $\Delta t = 4\times10^{-4}$--$2.5\times10^{-5}$,
\texttt{EPS} = $10^{-8}$); only the self-convergence error of $\psitd$ between
the two finest steps rises from $6$--$8\times10^{-6}$ to $1.5\times10^{-5}$, its
floor.  With \texttt{NANDER} = 0 the slow \texttt{LQ*8} mode still decides most
steps at \texttt{NRMAX} $\ge$ 240, and the floor term changes little.

\subsection{New defaults and V5.52 compatibility}

\begin{center}
\begin{tabular}{|l|c|c|l|}\hline
parameter & V5.52 & V5.60 & meaning \\\hline
\texttt{FSLP}    & 0.3 & 1.0 & $k_L$ \\
\texttt{MDLSOL}  & (0) & 1   & ramp width $\delta$ \\
\texttt{MDLNZD}  & (0) & 1   & bundle transport with the source-slaved C$^0$ part \\
\texttt{MDLWBC}  & (0) & 1   & wall density condition \\
\texttt{MDLN0FL} & (0) & 0   & Larsen limiter on $D_{01}$ (2: element form; 1: legacy nodal form) \\
\texttt{FLlim}   & --  & 0.25 & limiter factor $\alpha$ (C$^0$; $D_{01}$ with \texttt{MDLN0FL} = 1, 2) \\
\texttt{RLXN0}   & --  & 0.5 & neutral-only relaxation (legacy \texttt{MDLN0FL} = 1 only) \\
\texttt{NANDER}  & (0) & 0   & Anderson acceleration of the Picard sweeps \\\hline
\end{tabular}
\end{center}

A value in parentheses means that the switch did not exist in V5.52 and that
the value reproduces the V5.52 behaviour.  These are all the defaults that differ.
\texttt{parm/txparm.v552default} sets
\texttt{FSLP} = 0.3, \texttt{MDLSOL} = \texttt{MDLNZD} = \texttt{MDLWBC} =
\texttt{MDLN0FL} = 0.  With it V5.60 is not bit-identical to V5.52.20 (the
optimized build schedules the changed routines differently), but one step from
the same state gives double-precision states that agree to $10^{-12}$ (median
over the variables) and $10^{-9}$ at worst, except $\psitd$ at its Picard noise
floor ($3\times10^{-7}$); the profile output agrees to its single-precision
resolution.  The table \texttt{data/Carbon-S0F0\_20260921.txt} is read at start-up in
either case.

\subsection{Other changes since V5.51}
\label{sec:v552other}

\begin{itemize}
\item \texttt{MDLPCK} = 2 (default): block-tridiagonal solver with one step of
  iterative refinement; the Picard floor falls from about $10^{-5}$ to
  $5\times10^{-7}$ on the worst equation.
\item \texttt{oldmixa} (default 0.5): Picard under-relaxation applied only on
  time steps whose sweep is measured not to contract.
\item BDF2 keeps its history across \texttt{CONT}; one backward-Euler step is
  taken where there is none (cold start, \texttt{LOAD}).
\item Analytic axis limit of the Pfirsch--Schl\"uter Fourier factors,
  $f_m(\cdot,0) = 0$; $\hat q^2$ at the axis from the interior instead of $q(0)^2$.
\item $D_{0\r{z}}$ uses the impurity neutral speed $V_{0z}$.
\item \texttt{FSTPTM} = 3 (default): multi-ion composition channel for the
  turbulent particle flux.  With \texttt{FSTPTM} = 2 the impurity is not diffused on
  its own density gradient, and in long runs a grid-scale mode of $n_\r{z}$ (wavelength
  about three mesh intervals, sign-alternating) grows: the JT-60U case of
  Sec.~\ref{sec:v561val} stops at 2.4 s with $T_\r{z} < 0$ at $\rho \approx 0.43$.  The mode
  is tied to the mesh, not to a physical scale, and is not a physical impurity
  instability; it is the discrete form of an under-regularized closure, in which nothing
  damps short-wavelength perturbations of $n_\r{z}$.  The diffusion of each thermal ion on
  its own gradient in \texttt{FSTPTM} = 3 damps them at the rate $D k^2$; the same case then
  runs to 4 s with no grid-scale component above the curvature of the profile.  In the NBI cases $\nee$ and $u_{\r{i}\phi}$ change
  by less than 1.5\% (in the ohmic default case $u_{\r{i}\phi} \approx 0.5$ km/s changes by
  15\%), the core impurity density by 25--55\%.  \texttt{FSTPTM} = 2 remains
  available and must be given for every species (\texttt{FSTPTM} = 2, 2, 2).
\item Turbulent force of \texttt{FSTPTM} = 2, 3 in the wall row.  The pressure-gradient part
  of the electron-driven turbulent force is assembled through the diamagnetic flow
  ($n\,BV_1$, LQ*8) and $n\Phi'$.  LQe8 is the algebraic definition
  $n_\r{e}BV_{1\r{e}} = -(I/\psi')(p_\r{e}'/(Z_\r{e}e) + n_\r{e}\Phi')$, so in the continuum this is
  the direct $\partial_V p_\r{e}$ and the force is
  \texttt{FSVAHL}\,$\partial_V p_\r{e}$ + (1 $-$ \texttt{FSVAHL})\,$T_\r{e}\partial_V\nee$.  Discretely
  LQ*8 is solved with the SUPG weights $N_j + p h N_j'$ (\texttt{iSUPG8} = 1) and the force rows
  with the Galerkin $N_i$: at an interior node the SUPG moments of the two elements around it
  cancel to $O(h^2)$, but the wall node has the last element only, and there the two forms
  differ by the order of the LQe8 residual.  The particle flux at the wall is fixed by the
  wall row of LQ*4, so this gave a spurious wall flux (JT-60U case, \texttt{NRMAX} = 60: an
  electron inflow of $0.22\times10^{20}$ s$^{-1}$).  The wall row (test function of node
  \texttt{NRMAX}) now uses the direct-gradient form for every species (that of
  \texttt{FSTPTM} = 0); not with \texttt{MDFIXT} $\neq$ 0.  The removed artificial part of the
  electron wall flux is 0.38 and 0.20 $\times10^{20}$ s$^{-1}$ at \texttt{NRMAX} = 60 and 120 (1.3 s) and 0.16, 0.10 and 0.045 at 60, 120 and 240 (0.05 s), so it goes with the mesh; the
  corrected wall flux itself converges slowly (+0.16 and +0.05 at 60 and 120; +0.18, +0.03 and +0.02 at 60, 120 and 240), being a small
  difference of large terms in the wall row.  The non-ambipolar part of the wall flux (the
  fast-ion return current), $j_r$, the $j\times B$ torque and $L_\phi$ are unchanged to three
  digits, the profiles change by less than 0.5\% (JT-60U case to 4 s, \texttt{FSRP} = 0 and 1,
  \texttt{FSTPTM} = 2 and 3, X-point and circular test equilibria), and the Picard counts are
  unchanged.
\item Recycling at the main wall.  The radial ion outflow through the wall $\rho_b$,
  $\Gamma_s(\rho_b)^+ = \max(n_s\aves{\vec{u}_s\cdot\nabla V}, 0)$ at \texttt{NRMAX}, returns as a
  wall-side inflow of neutrals like the gas puff: \texttt{RWALL} (default 0.98) of the bulk ions
  to $n_{01}$, \texttt{RWALLZ} (0.42) of the impurity ions to the impurity neutrals, and
  \texttt{RWALLIZ} (0, off) of the bulk ions to the impurity neutrals by sputtering; the impurity
  never returns as a hydrogen neutral.  This path is separate from the target recycling of the
  SOL parallel loss (\texttt{rGamm0}, \texttt{rGamm0zz}, \texttt{rGamm0iz}), which
  \texttt{tx\_sol\_particle\_balance} computes from a divertor recycling loop (escape to the
  upstream SOL against the pump, $c_1/(c_1+c_5)$ = 0.980 for the bulk ions; for the impurity the
  reflection and self-sputtering yield 0.42 times an escape factor 0.051, 0.0213; sputtering by
  the bulk ions 0.0099).  There is no such loop at the main wall, so \texttt{RWALL} is that of a
  saturated wall with a small retention and \texttt{RWALLZ} the yield itself.  In the JT-60U case
  the main wall returns $0.40\times10^{20}$ s$^{-1}$ (target recycling 14.4, gas puff 28.1); $\nee$
  changes by 0.2\% (0.6\% with \texttt{FSRP} = 1), and even \texttt{RWALLIZ} = 0.15 changes the core
  impurity density by 0.04\%, the wall source being screened by the SOL.  \texttt{W} (TXSTAT) and
  \texttt{O 7} (\texttt{tx\_sources.dat}) print the SOL parallel loss, the wall fluxes and the
  neutral sources by path (target, main wall, gas puff) with the ionization.  The three
  coefficients are kept in the save file in a second extension record; older files load with
  the current values and newer files load in older builds.
\item Finite-collisionality closure for the higher-order flows; triangularity of
  grid-traced equilibria regularized; G-EQDSK $p'$ normalization corrected;
  \texttt{ipbtdir} = 1 by default.
\item \texttt{MDLNSUP} (default 1): the SOL upstream densities of the ions and the
  impurity, \texttt{PNsSUP(2:3)}, are derived from \texttt{PNsSUP(1)} and
  \texttt{Zeffin} (quasi-neutral, $Z_\r{eff}$ = \texttt{Zeffin}) whenever the
  parameters are read; before, they were set once from the defaults and a namelist
  value of \texttt{PNsSUP(1)} or \texttt{Zeffin} left them unchanged.
  \texttt{MDLNSUP} = 0 uses \texttt{PNsSUP(1:3)} as given.  The SOL loss
  $-\nu_L(n_s - n_{\r{SUP},s})$ holds the SOL density near $n_{\r{SUP},s}$, so
  \texttt{PNsSUP(1)} acts as the upstream SOL density and sets the density level of
  the whole plasma.
\item Diagnostic switches, off by default: \texttt{iaxproj} (near-axis
  regularity projection), \texttt{Dnfac} (particle diffusion in the impurity
  continuity equation).
\end{itemize}

\begin{thebibliography}{9}
\bibitem{NF2008} M. Honda, T. Takizuka, A. Fukuyama, M. Yoshida and T. Ozeki,
  Nucl. Fusion \textbf{48} (2008) 085003.
\bibitem{Yoshida2006} M. Yoshida et al., Plasma Phys. Control. Fusion \textbf{48} (2006) 1673.
\bibitem{KA1995} M. Kikuchi and M. Azumi, Plasma Phys. Control. Fusion \textbf{37} (1995) 1215,
  Appendix.
\bibitem{Tala2012} T. Tala et al., Proc. 24th IAEA Fusion Energy Conf. (San Diego, 2012) ITR/P1-19.
\bibitem{Tala2011} T. Tala et al., Nucl. Fusion \textbf{51} (2011) 123002.
\bibitem{Mantica2010} P. Mantica et al., Phys. Plasmas \textbf{17} (2010) 092505.
\bibitem{Weisen2012} H. Weisen et al., Nucl. Fusion \textbf{52} (2012) 042001; \textbf{52} (2012) 114024.
\bibitem{Honda2013} M. Honda et al., Nucl. Fusion \textbf{53} (2013) 073050.
\bibitem{Casiraghi2021} I. Casiraghi et al., Nucl. Fusion \textbf{61} (2021) 116068.
\end{thebibliography}

\end{document}
